% =====================================================================
%  LUCRARE DE DIPLOMĂ — Universitatea Politehnica Timișoara
%  Facultatea de Automatică și Calculatoare — Informatică Zi
%  Compilare:  latexmk -xelatex thesis.tex
% =====================================================================
\documentclass[a4paper,12pt]{scrreprt}

% ── Căi ──────────────────────────────────────────────────────────────
\def \cls {./cls}
\def \pics {./pics}
\def \chapters {./chapters}

% ── Clasa UPT (limba română, licență) + bibliografie + glosar ────────
\usepackage[language=romanian,type=bachelor]{\cls/standard}
\addbibresource{\cls/bib.bib}
% ── Glosar de acronime — Sistem Smart Home Local-First ──────────────
\newacronym{iot}{IoT}{Internet of Things}
\newacronym{ml}{ML}{Machine Learning}
\newacronym{if}{IF}{Isolation Forest}
\newacronym{mqtt}{MQTT}{Message Queuing Telemetry Transport}
\newacronym{api}{API}{Application Programming Interface}
\newacronym{rest}{REST}{Representational State Transfer}
\newacronym{jwt}{JWT}{JSON Web Token}
\newacronym{json}{JSON}{JavaScript Object Notation}
\newacronym{http}{HTTP}{HyperText Transfer Protocol}
\newacronym{tls}{TLS}{Transport Layer Security}
\newacronym{mjpeg}{MJPEG}{Motion JPEG}
\newacronym{mdns}{mDNS}{Multicast Domain Name System}
\newacronym{lwt}{LWT}{Last Will and Testament}
\newacronym{nvs}{NVS}{Non-Volatile Storage}
\newacronym{gpio}{GPIO}{General-Purpose Input/Output}
\newacronym{adc}{ADC}{Analog-to-Digital Converter}
\newacronym{pwm}{PWM}{Pulse-Width Modulation}
\newacronym{i2c}{I2C}{Inter-Integrated Circuit}
\newacronym{spi}{SPI}{Serial Peripheral Interface}
\newacronym{uart}{UART}{Universal Asynchronous Receiver-Transmitter}
\newacronym{pir}{PIR}{Passive Infrared}
\newacronym{ldr}{LDR}{Light-Dependent Resistor}
\newacronym{led}{LED}{Light-Emitting Diode}
\newacronym{dht}{DHT}{Digital Humidity and Temperature}
\newacronym{rpi}{RPi}{Raspberry Pi}
\newacronym{cpu}{CPU}{Central Processing Unit}
\newacronym{ram}{RAM}{Random-Access Memory}
\newacronym{sram}{SRAM}{Static Random-Access Memory}
\newacronym{psram}{PSRAM}{Pseudo-Static Random-Access Memory}
\newacronym{wifi}{Wi-Fi}{Wireless Fidelity}
\newacronym{lan}{LAN}{Local Area Network}
\newacronym{ip}{IP}{Internet Protocol}
\newacronym{tcp}{TCP}{Transmission Control Protocol}
\newacronym{sql}{SQL}{Structured Query Language}
\newacronym{orm}{ORM}{Object-Relational Mapping}
\newacronym{ui}{UI}{User Interface}
\newacronym{ux}{UX}{User Experience}
\newacronym{gdpr}{GDPR}{General Data Protection Regulation}
\newacronym{auc}{AUC}{Area Under the Curve}
\newacronym{roc}{ROC}{Receiver Operating Characteristic}
\newacronym{qos}{QoS}{Quality of Service}
\newacronym{esp32}{ESP32}{Espressif Systems 32-bit SoC}
\newacronym{oasis}{OASIS}{Organization for the Advancement of Structured Information Standards}
\newacronym{svm}{SVM}{Support Vector Machine}
\newglossaryentry{edge}{name={edge},description={paradigmă de procesare a datelor la marginea rețelei, în apropierea sursei de date}}


% ── Așezare în pagină ────────────────────────────────────────────────
\setlength{\parindent}{1.5cm}
\usepackage{setspace}
\onehalfspacing
\usepackage{graphicx}
\usepackage{caption}
\usepackage{booktabs}
\usepackage{float}      % opțiunea [H] pentru figuri/tabele
\usepackage{textcomp}   % simboluri text suplimentare
\usepackage[htt]{hyphenat} % permite despărțirea în silabe a textului \texttt (cod în text)

% Macro-uri utilizate în capitole, dar nedefinite de clasă
\providecommand{\listingref}[1]{Listarea~\ref{#1}}
\providecommand{\textOmega}{\ensuremath{\Omega}}
\newcommand{\slb}{/\allowbreak}   % slash cu punct de rupere pentru topicuri/căi lungi

% ── Robustețe la compilare (cerute explicit) ─────────────────────────
\usepackage{microtype}
\usepackage{xurl}
\setlength{\emergencystretch}{3em}
\setlength{\headheight}{15pt}

% ── Listinguri de cod: TypeScript (nu e inclus implicit) + C++ ───────
\lstdefinelanguage{TypeScript}{
  morekeywords={abstract,as,async,await,break,case,catch,class,const,%
    continue,debugger,default,delete,do,else,enum,export,extends,false,%
    finally,for,from,function,get,if,implements,import,in,instanceof,%
    interface,let,new,null,of,private,protected,public,readonly,return,%
    set,static,super,switch,this,throw,true,try,type,typeof,var,void,%
    while,yield},
  morekeywords=[2]{string,number,boolean,any,unknown,Promise,Record,Array,%
    Partial,void},
  sensitive=true,
  morecomment=[l]{//},
  morecomment=[s]{/*}{*/},
  morestring=[b]',
  morestring=[b]",
  morestring=[b]`
}
\lstset{
  basicstyle=\ttfamily\scriptsize,
  columns=fullflexible,
  breaklines=true,
  breakatwhitespace=false,
  keepspaces=true,
  showstringspaces=false,
  tabsize=2,
  frame=single,
  captionpos=b
}

% ── Pagina de titlu cu logo-uri (UPT stânga, AC dreapta) ─────────────
\makeatletter
\renewcommand{\maketitle}{
\begin{titlepage}
\noindent

\includegraphics[height=2.1cm,keepaspectratio]{\pics/logo/upt}%
\hfill

\includegraphics[height=2.1cm,keepaspectratio]{\pics/logo/ac}
\vspace{1em}
\begin{center}
Universitatea Politehnica Timișoara\\
Facultatea de Automatică și Calculatoare\\
Ingineria Sistemelor\\
\vspace{5em}
{\Huge\textbf{LUCRARE DE DIPLOMĂ}}
\vspace{2.5em}

{\huge\textbf{\@title}}
\end{center}
\vspace{3.5em}
\begin{flushright}
\Large{Absolvent\\
\@author}
\end{flushright}
\vspace{2em}
\Large{
\begin{flushleft}
Coordonator științific\\
Ș.l. dr. ing. Paul NEGÎRLA\\[0.9em]
Îndrumător\\
Asist. univ. drd. ing. Darius NEGÎRLA
\end{flushleft}
\vfill
\begin{center}
\Large{Timișoara, \the\year}
\end{center}
}
\end{titlepage}
}
\makeatother

% =====================================================================
\begin{document}
% ── Front matter (numerotare romană) ─────────────────────────────────
\pagenumbering{roman}

\title{Sistem Smart Home Local-First cu Detecția Anomaliilor Comportamentale prin Învățare Automată pe Edge}
\author{Vidac Denisa}
\date{2026}

\maketitle

% ── Rezumat (RO) ─────────────────────────────────────────────────────
\renewcommand{\abstractname}{Rezumat}
\begin{abstract}
\noindent
Lucrarea prezintă proiectarea și implementarea unui sistem \textit{smart home}
\textit{local-first}, care funcționează integral în rețeaua locală a locuinței,
fără dependență de servicii cloud. Sistemul este organizat pe cinci straturi:
senzori și actuatori, un nod ESP32 cu firmware Arduino, un \gls{rpi} care găzduiește
brokerul \gls{mqtt} \textit{Mosquitto}, o bază de date PostgreSQL/TimescaleDB, un
API NestJS și o aplicație mobilă React~Native. Nodul interior achiziționează
temperatura și umiditatea (două senzoare \gls{dht}11), nivelul de gaz (MQ-2) și
luminozitatea (două \gls{ldr}), și comandă trei \gls{led}-uri, un ventilator, un
servomotor și un buzzer; o cameră ESP32-CAM cu senzor \gls{pir} acoperă exteriorul.
Contribuția principală o reprezintă un modul de \textbf{detecție de anomalii
comportamentale} bazat pe algoritmul \gls{if} (\textit{Isolation Forest}), antrenat
și executat \textbf{local}, pe Raspberry Pi, care învață tiparele normale ale
locuinței și semnalează abaterile fără reguli prestabilite. Lucrarea descrie
arhitectura, implementarea prototipului și metodologia de evaluare, evidențiind
avantajele de confidențialitate, cost și autonomie ale abordării locale.

\medskip\noindent\textbf{Cuvinte-cheie:} smart home, Internet of Things, local-first,
detecție de anomalii, Isolation Forest, edge computing, ESP32, MQTT, Raspberry Pi.
\end{abstract}

% ── Abstract (EN) ────────────────────────────────────────────────────
\renewcommand{\abstractname}{Abstract}
\begin{abstract}
\noindent
This thesis presents the design and implementation of a \textit{local-first}
smart home system that operates entirely within the home's local network, with no
dependency on cloud services. The system is organized in five layers: sensors and
actuators, an ESP32 node running Arduino firmware, a Raspberry~Pi hosting the
\textit{Mosquitto} \gls{mqtt} broker, a PostgreSQL/TimescaleDB database, a NestJS
API and a React~Native mobile application. The indoor node acquires temperature and
humidity (two \gls{dht}11 sensors), gas level (MQ-2) and illumination (two
\gls{ldr}s), and drives three \gls{led}s, a fan, a servo motor and a buzzer; an
ESP32-CAM with a \gls{pir} sensor covers the outdoor area. The main contribution is
a \textbf{behavioural anomaly detection} module based on the \gls{if}
(\textit{Isolation Forest}) algorithm, trained and executed \textbf{locally} on the
Raspberry~Pi, which learns the household's normal patterns and flags deviations
without predefined rules. The thesis describes the architecture, the prototype
implementation and the evaluation methodology, highlighting the privacy, cost and
autonomy benefits of the local approach.

\medskip\noindent\textbf{Keywords:} smart home, Internet of Things, local-first,
anomaly detection, Isolation Forest, edge computing, ESP32, MQTT, Raspberry Pi.
\end{abstract}

% ── Cuprins și liste ─────────────────────────────────────────────────
\tableofcontents
\listoffigures
\addcontentsline{toc}{chapter}{\listfigurename}
\listoftables
\addcontentsline{toc}{chapter}{\listtablename}

\printnoidxglossaries
\clearpage

% ── Main matter (numerotare arabă) ───────────────────────────────────
\pagenumbering{arabic}
\pagestyle{headings}
\markboth{Universitatea Politehnica Timișoara \\ Informatică Zi, 2026 \\ Vidac Denisa}{Sistem Smart Home Local-First}

% cap1_introducere.tex
\chapter{Introducere}
\label{ch:introducere}

% ─────────────────────────────────────────────────────────────────────────────
\section{Context și motivație}
\label{sec:motivatie}

Locuințele inteligente reprezintă una dintre aplicațiile cu cea mai rapidă
expansiune din domeniul \gls{iot}. Dacă la începutul deceniului trecut
termenul „smart home" evoca un lux rezervat câtorva entuziaști, astăzi
milioane de gospodării integrează senzori de temperatură, camere de
supraveghere, sisteme de iluminat și asistenți vocali conectați la internet.
Piața globală a locuințelor inteligente a depășit \$80 de miliarde în 2023
și se preconizează o creștere susținută în anii următori
\cite{stojkoska2017smarthome,gubbi2013iot}. La baza acestei transformări stă
tocmai convergența \gls{iot}: dispozitive cu resurse limitate care colectează
date din mediul fizic, le transmit în rețea și declanșează acțiuni automatizate
în funcție de starea mediului \cite{atzori2010iot}.

Cu toate acestea, marea majoritate a soluțiilor comerciale disponibile astăzi
(Google Home, Amazon Alexa, Apple HomeKit, Tuya etc.) se bazează pe un model
\emph{cloud-first}: datele brute ale senzorilor sunt trimise pe servere aflate
în altă țară, procesate acolo și returnate aplicației mobile printr-un lanț
de apeluri de rețea. Această arhitectură ridică o serie de probleme grave:

\begin{itemize}
  \item \textbf{Confidențialitate și protecția datelor.}
  Informații sensibile privind comportamentul cotidian al locatarilor
  (ore de prezență acasă, temperaturi, mișcări în cameră) sunt stocate
  pe infrastructuri terțe, adesea în afara jurisdicției \gls{gdpr}
  \cite{zheng2018smarthomeprivacy,gdpr2016}.
  Chiar atunci când furnizorii respectă legislația europeană, utilizatorul
  nu deține controlul efectiv asupra datelor generate în propria locuință
  \cite{alrawi2019sokhome}.

  \item \textbf{Dependența de conectivitatea la internet.}
  Dacă furnizorul de servicii întrerupe suportul pentru un produs sau dacă
  conexiunea la internet este indisponibilă temporar, întregul sistem poate
  deveni nefuncțional. Experiențele legate de retragerea prematură a unor
  platforme (Google Stadia, SmartThings Hub v2 etc.) au demonstrat că
  \emph{vendor lock-in}-ul reprezintă un risc real pentru utilizatorii care
  investesc în ecosisteme proprietare \cite{kleppmann2019localfirst}.

  \item \textbf{Cost și complexitate.}
  Abonamentele lunare pentru platformele cloud, costul gateway-urilor
  proprietare și imposibilitatea de a integra hardware eterogen cresc
  semnificativ bariera de acces. Soluțiile open-source existente (Home
  Assistant, OpenHAB) oferă alternativă, dar cer cunoștințe tehnice
  avansate și nu includ componente de analiză inteligentă a comportamentului.

  \item \textbf{Latență și procesare la margine rețelei.}
  Trimiterea fiecărei citiri de senzor spre cloud și așteptarea unui răspuns
  introduce latențe inacceptabile pentru scenarii de siguranță (detectare gaz,
  alertă de intruziune nocturnă). Paradigma \emph{edge computing} propune
  procesarea locală a datelor, acolo unde sunt generate
  \cite{shi2016edge}.
\end{itemize}

Aceste limitări nu sunt de natură exclusiv tehnică: ele reflectă un dezechilibru
structural în care controlul locuinței proprii este delegat unor entități externe.
Lucrarea de față propune o arhitectură alternativă, \emph{local-first}, care
rezolvă fiecare dintre problemele enumerate anterior.

% ─────────────────────────────────────────────────────────────────────────────
\section{Problema abordată și abordarea local-first}
\label{sec:problema}

Problema centrală abordată în această lucrare este următoarea:
\emph{Este posibil să se construiască un sistem de locuință inteligentă complet
funcțional, capabil de monitorizare în timp real, control al actuatorilor,
supraveghere video și detecție inteligentă a anomaliilor comportamentale,
folosind exclusiv infrastructură locală, fără nicio dependență de servicii cloud,
și cu un cost total de hardware accesibil?}

Răspunsul propus este o arhitectură \textbf{local-first}
\cite{kleppmann2019localfirst}: toate datele rămân în rețeaua locală (\gls{lan})
a locatarului, procesarea se realizează pe un gateway local (\gls{rpi}), iar
aplicația mobilă comunică direct cu serverul intern fără a tranzita internetul.
Această abordare se înscrie în tendința mai largă de \emph{edge computing}
\cite{shi2016edge,premsankar2018edgeiot}: în loc să se centralizeze inteligența
în cloud, logica de decizie este mutată cât mai aproape de sursele de date
(senzori, camere, actuatori).

Sistemul implementat funcționează integral în \gls{lan}, inclusiv atunci când
conexiunea la internet este indisponibilă. Singurul trafic care poate opțional
ieși din rețeaua locală este notificarea push, iar și aceasta poate fi dezactivată.
Datele personale ale utilizatorului (temperaturi, ore de prezență, imagini de la
cameră) nu părăsesc niciodată locuința fără consimțământ explicit.

Componenta de \gls{ml} — motorul de detecție a anomaliilor comportamentale bazat
pe algoritmul \gls{if} — rulează direct pe \gls{rpi}, eliminând necesitatea
unui serviciu cloud de inferență. Aceasta este contribuția originală centrală a
lucrării, detaliată în \secref{sec:contributii} și dezvoltată extensiv în
capitolele \chapref{ch:metodologie} și \chapref{ch:implementare}.

% ─────────────────────────────────────────────────────────────────────────────
\section{Obiectivele lucrării}
\label{sec:obiective}

Lucrarea își propune proiectarea, implementarea și evaluarea unui sistem de
locuință inteligentă local-first, cu scopuri concrete structurate pe cinci
straturi arhitecturale. Obiectivele principale sunt:

\begin{enumerate}
  \item \textbf{Sistem funcțional pe cinci straturi.}
  Proiectarea unei arhitecturi coerente care integrează: senzori și actuatori
  fizici, noduri \gls{iot} bazate pe microcontrolere ESP32 programate în
  Arduino C++, un gateway local pe \gls{rpi}, un \gls{api} \gls{rest} +
  WebSocket și o aplicație mobilă — toate comunicând în \gls{lan} fără cloud.

  \item \textbf{Achiziție de date cu hardware real.}
  Conectarea și programarea celor 11 componente hardware reale ale nodului
  interior (ESP32 WROOM-32 cu două senzoare \gls{dht}11, un senzor de gaz
  MQ-2, două rezistențe fotosensibile \gls{ldr}, trei \gls{led}-uri, un
  buzzer, un servomotor SG90 și un ventilator 5V) și a nodului exterior
  (ESP32-CAM cu senzor \gls{pir} și cameră OV2640).

  \item \textbf{Control al actuatorilor în timp real.}
  Implementarea unui protocol de comandă bidirecțional prin \gls{mqtt} care
  permite aprinderea/stingerea \gls{led}-urilor, pornirea ventilatorului,
  mișcarea servo-ului (draperii), activarea buzzerului și configurarea \gls{wifi}
  — toate inițiate din aplicația mobilă cu latență sub o secundă în \gls{lan}.

  \item \textbf{Motor de \gls{ml} local pentru detecția anomaliilor.}
  Implementarea unui pipeline complet de \gls{if} (antrenare, normalizare,
  inferență în timp real, reantrenare automată zilnică) care rulează exclusiv
  pe \gls{rpi}, fără apeluri externe, și care generează alerte de tip
  \texttt{ML\_ANOMALY} integrate în sistemul de notificări al aplicației mobile.

  \item \textbf{Aplicație mobilă cross-platform.}
  Dezvoltarea unei aplicații React Native / Expo cu patru tab-uri funcționale
  (Dashboard cu date live, Istoric cu grafice, Alerte, Control actuatori +
  cameră), care se conectează automat la gateway prin \gls{mdns} sau adresă IP
  și oferă autentificare \gls{jwt}.

  \item \textbf{Independență completă de cloud.}
  Demonstrarea că întregul sistem (monitorizare, control, \gls{ml}, video)
  funcționează fără nicio dependență de servicii externe, asigurând
  confidențialitatea datelor și reziliența la lipsa conectivității la internet.

  \item \textbf{Cost redus și reproductibilitate.}
  Construirea sistemului din componente disponibile comercial, cu un cost total
  de hardware accesibil, și documentarea completă a firmware-ului, configurației
  serverului și structurii bazei de date pentru a permite reproducerea independentă.
\end{enumerate}

% ─────────────────────────────────────────────────────────────────────────────
\section{Contribuții originale}
\label{sec:contributii}

Principala contribuție originală a lucrării este \textbf{implementarea unui motor
de detecție a anomaliilor comportamentale bazat pe algoritmul
\gls{if} \cite{liu2008isolation}, rulat în întregime local pe \gls{rpi}}.

Algoritmul \gls{if}, propus de Liu, Ting și Zhou \cite{liu2008isolation}, izolează
punctele anomale printr-un mecanism de partiționare aleatoare a spațiului de
caracteristici: punctele rare (anomale) sunt separate prin mai puțini pași
față de punctele normale, rezultând o \emph{lungime de cale medie} mai mică și,
implicit, un scor de anomalie mai mare. Această proprietate îl face deosebit de
potrivit pentru scenarii de tip \emph{edge}: nu necesită etichete de antrenare
(algoritm nesupervizat), are complexitate computațională redusă și poate fi
reantrenant periodic pe date noi fără intervenție umană.

În contextul acestui sistem, motorul \gls{ml} lucrează cu 12 trăsături extrase
din citirile senzorilor (temperaturi de la ambele senzoare DHT11, umiditate,
nivel de gaz normalizat, intensitate luminoasă normalizată logaritmic, prezența
mișcării, ora zilei, ziua săptămânii, indicator de weekend/noapte, variația
temperaturii pe o oră și numărul de evenimente de mișcare în 10 minute).
Modelul este antrenat pe ultimele 30 de zile de date și reantrenant automat
zilnic la ora 03:00 printr-un serviciu systemd. Inferența se realizează în timp
real la fiecare citire nouă de senzori, fără latența unui apel de rețea spre
cloud.

Această abordare reprezintă un avans față de soluțiile comerciale existente, care
fie nu includ deloc componente de \gls{ml}, fie externalizează procesarea spre
servicii cloud cu implicații de confidențialitate. Spre deosebire de metodele
supervizate care necesită date etichetate (greu de obținut în medii rezidențiale),
\gls{if} este \emph{nesupervizat} și se adaptează automat la rutina specifică a
locatarilor \cite{chandola2009anomaly,fahim2019anomaly,cook2020anomalyiot}.

Contribuțiile secundare ale lucrării includ:
\begin{itemize}
  \item \textbf{Integrarea completă firmware↔broker↔API↔aplicație mobilă.}
  Lanțul complet de la firmware Arduino C++ pe ESP32, prin brokerul Mosquitto
  \gls{mqtt}, prin \gls{api}-ul NestJS cu TypeORM și TimescaleDB, până la
  interfața React Native — implementat și funcțional pe hardware fizic.

  \item \textbf{Arhitectura de date time-series.}
  Utilizarea TimescaleDB (extensie PostgreSQL orientată pe serii de timp) pentru
  stocarea eficientă a citirilor de senzori, cu indexare automată pe timestamp
  și suport pentru interogări de agregare pe intervale temporale.

  \item \textbf{Provisioning Wi-Fi și descoperire automată în rețea.}
  Mecanism de configurare inițială prin hotspot propriu al nodului ESP32
  (\texttt{SmartHome-Setup}), cu stocare a credențialelor în \gls{nvs} și
  descoperire automată a gateway-ului prin \gls{mdns} (\texttt{smarthome.local}).

  \item \textbf{Modul de simulare pentru dezvoltare fără hardware.}
  Un modul \texttt{simulator/} integrat în backend-ul NestJS care generează date
  realiste ale senzorilor (inclusiv variabilele extrase ulterior de motorul \gls{ml}),
  permițând dezvoltarea și testarea completă a stivei software fără acces la
  dispozitivele fizice.
\end{itemize}

% ─────────────────────────────────────────────────────────────────────────────
\section{Structura lucrării}
\label{sec:structura}

Lucrarea este organizată în șase capitole, fiecare abordând un aspect distinct
al sistemului:

\textbf{\chapref{ch:introducere} — Introducere} (capitolul de față) prezintă
contextul domeniului, motivațiile care au condus la abordarea local-first,
obiectivele concrete ale sistemului și contribuțiile originale.

\textbf{\chapref{ch:studii} — Studiu bibliografic} trece în revistă literatura
relevantă: evoluția sistemelor \gls{iot} și a locuințelor inteligente, protocoale
de comunicație cu accent pe \gls{mqtt}, tehnici de detecție a anomaliilor în date
de senzori (comparație între \gls{if}, LOF, One-Class SVM și metode bazate pe
rețele neuronale), paradigma \emph{edge computing} și cadrul legal \gls{gdpr}
privind datele personale în medii rezidențiale.

\textbf{\chapref{ch:metodologie} — Metodologie și proiectare} descrie arhitectura
pe cinci straturi a sistemului, deciziile de proiectare (alegerea componentelor
hardware, a stivei software, a bazei de date time-series), modelarea entităților,
fluxurile de date \gls{mqtt} și pipeline-ul complet al motorului \gls{ml}.

\textbf{\chapref{ch:implementare} — Implementare} detaliază realizarea efectivă a
fiecărui strat: firmware-ul Arduino C++ al nodului interior și al ESP32-CAM,
configurarea brokerului Mosquitto și a \gls{rpi} ca gateway, implementarea
\gls{api}-ului NestJS cu modulele de \gls{mqtt}, WebSocket, autentificare \gls{jwt}
și \gls{ml}, și dezvoltarea aplicației mobile React Native cu toate funcționalitățile
sale.

\textbf{\chapref{ch:rezultate} — Rezultate și evaluare} prezintă rezultatele
obținute: funcționarea sistemului pe hardware real, evaluarea motorului de \gls{ml}
pe date ilustrative (cu date preliminare marcate explicit ca atare), latențele
măsurate în \gls{lan} și o discuție critică a limitărilor actuale și a direcțiilor
de îmbunătățire.

\textbf{\chapref{ch:concluzii} — Concluzii} sintetizează contribuțiile lucrării,
evaluează măsura în care obiectivele au fost atinse și propune direcții de
cercetare și dezvoltare ulterioară, inclusiv extinderea la mai multe noduri,
integrarea unor senzori suplimentari și posibilitatea unui model \gls{ml} mai
sofisticat (autoencodere, LSTM) pentru detecția anomaliilor.

% cap2_studii.tex
\chapter{Studii anterioare}
\label{ch:studii}

Prezentul capitol trece în revistă stadiul actual al domeniilor relevante pentru această
lucrare: platformele de automatizare rezidențială, protocoalele de comunicație \gls{iot},
paradigmele de calcul la margine (edge/fog) și metodele de detecție a anomaliilor comportamentale.
Pe baza acestei analize, sunt identificate lacunele din literatura de specialitate pe care
sistemul propus urmărește să le acopere.

% ======================================================================
\section{Platforme smart home comerciale și open-source}
\label{sec:studii-platforme}
% ======================================================================

Piața automatizărilor rezidențiale a cunoscut o expansiune rapidă în ultimul deceniu, atât
în segmentul comercial cât și în cel open-source. Stojkoska și Trivodaliev \cite{stojkoska2017smarthome}
oferă o trecere în revistă cuprinzătoare a provocărilor și soluțiilor din domeniu, evidențiind
că interoperabilitatea, securitatea și confidențialitatea datelor rămân problemele centrale
nerezolvate pe deplin de nicio platformă.

\subsection{Platforme comerciale}

Printre soluțiile comerciale cu cea mai largă adopție se numără Amazon Alexa/Echo,
Google Home și Apple HomeKit. Toate trei se bazează pe arhitecturi hibride cloud–dispozitiv:
comenzile utilizatorului sunt transmise serverelor producătorului, procesate în cloud și
returnate dispozitivului. Această abordare oferă avantaje clare — actualizări automate,
recunoaștere vocală de înaltă calitate, ecosisteme bogate de aplicații — dar implică o
serie de dezavantaje fundamentale. Alam et al.\ \cite{alam2012smarthome} subliniază că
dependența de conectivitate la internet face ca astfel de sisteme să devină nefuncționale
în absența conexiunii sau la întreruperea serviciului cloud al producătorului.

Aspectul confidențialității este tratat în detaliu de Zheng et al.\ \cite{zheng2018smarthomeprivacy},
care arată că datele colectate de dispozitivele smart home comerciale — inclusiv obiceiuri
zilnice, prezența la domiciliu sau conversații ambientale — sunt transferate și stocate pe
servere terțe, deseori în jurisdicții cu standarde de protecție a datelor diferite. O
analiză de securitate sistematică realizată de Alrawi et al.\ \cite{alrawi2019sokhome} pe
un eșantion larg de dispozitive IoT de consum a identificat vulnerabilități la nivelul
firmware-ului, al protocoalelor de comunicație și al \gls{api}-urilor cloud, concluzionând
că nivelul de securitate al ecosistemelor comerciale rămâne nesatisfăcător în raport cu
amenințările actuale.

\subsection{Platforme open-source}

La polul opus se situează platformele open-source, dintre care Home Assistant și openHAB
sunt cele mai reprezentative. Home Assistant, scris în Python și distribuit sub licență
Apache 2.0, rulează local pe hardware modest (inclusiv Raspberry Pi) și oferă o interfață
web completă, suport pentru sute de integrări și un motor de automatizare bazat pe
evenimente. openHAB, scris în Java, adoptă o filosofie similară de independență față de
cloud, cu accent pe extensibilitate prin plugin-uri.

Avantajul esențial al acestor platforme este că întreaga logică de control și toate datele
rămân în rețeaua locală a utilizatorului. Totuși, curba de învățare este semnificativă,
iar suportul pentru \gls{ml} nativ (mai ales pentru detecția de anomalii) este limitat sau
realizat prin integrări externe.

Sistemul propus în această lucrare împărtășește filosofia local-first a platformelor
open-source, dar se distinge prin integrarea unui motor de \gls{ml} propriu — Isolation
Forest — rulat direct pe gateway-ul Raspberry Pi, fără a depinde de nicio bibliotecă
externă de automatizare generică.

% ======================================================================
\section{Protocoale de comunicație IoT}
\label{sec:studii-protocoale}
% ======================================================================

Alegerea protocolului de comunicație are implicații directe asupra consumului de resurse,
a latenței, a scalabilității și a fiabilității unui sistem \gls{iot}. Cele mai utilizate
protocoale în domeniu sunt \gls{mqtt}, CoAP și HTTP. Naik \cite{naik2017mqttcomparison}
realizează o comparație detaliată a acestora, evidențiind că nu există o soluție universală
optimă, alegerea depinzând de constrângerile concrete ale aplicației.

\subsection{Comparație MQTT / CoAP / HTTP}

Tabelul~\ref{tab:protocoale} sintetizează principalele caracteristici ale celor trei
protocoale din perspectiva aplicabilității \gls{iot}.

\begin{table}[H]
\centering
\caption{Comparație MQTT, CoAP și HTTP pentru aplicații IoT}
\label{tab:protocoale}
\begin{tabular}{@{}p{3.2cm} p{3.8cm} p{3.8cm} p{3.8cm}@{}}
\toprule
\textbf{Caracteristică} & \textbf{MQTT} & \textbf{CoAP} & \textbf{HTTP/1.1} \\
\midrule
Transport            & TCP           & UDP           & TCP \\
\addlinespace
Model comunicație    & Publish/Subscribe (broker central) & Request/Response (P2P sau proxy) & Request/Response (client–server) \\
\addlinespace
Overhead antet       & Minim (2–5 B antet fix) & Redus (4 B antet fix) & Mare (text, zeci–sute de octeți) \\
\addlinespace
QoS                  & 3 niveluri (0, 1, 2) & 2 tipuri (confirmabil / neconfirmabil) & Implicit (TCP garantează livrarea) \\
\addlinespace
Last Will / Retained & Da (suport nativ) & Nu (necesită extensii) & Nu \\
\addlinespace
Potrivire IoT        & Excelentă (multe-la-mulți, noduri cu resurse limitate) & Bună (rețele cu pierderi, IPv6/6LoWPAN) & Moderată (overhead ridicat, nepotrivit pentru microcontrolere) \\
\addlinespace
Securitate           & TLS/SSL opțional & DTLS opțional & TLS/SSL standard \\
\addlinespace
Standard             & OASIS \cite{oasis2019mqtt} & IETF RFC 7252 \cite{rfc7252} & IETF RFC 7230 \\
\bottomrule
\end{tabular}
\end{table}

\subsection{Justificarea alegerii MQTT}

\gls{mqtt} a fost conceput explicit pentru conexiuni cu lățime de bandă mică și dispozitive
cu resurse limitate \cite{oasis2019mqtt}. Modelul publish/subscribe decuplează producătorul
de date (nodul ESP32) de consumatori (backend NestJS, motor ML), permitând mai multor
abonați să recepționeze simultan aceleași date fără ca nodul să cunoască numărul sau
identitatea acestora.

Caracteristica \emph{Last Will and Testament} este esențială pentru monitorizarea stării
nodurilor: la conectare, fiecare nod înregistrează un mesaj de „ultimă voință" pe topicul
\texttt{smarthome/\{nodeId\}/status}; dacă conexiunea se întrerupe neașteptat, brokerul
publică automat acest mesaj, informând restul sistemului. Mesajele \emph{retained} asigură
că un abonat nou primește imediat ultima valoare a unui topic, fără a aștepta o nouă
publicare.

Mosquitto, implementarea de referință pentru medii embedded \cite{light2017mosquitto},
rulează eficient pe Raspberry Pi 3B+ și suportă autentificare prin nume de utilizator și
parolă, acoperind cerințele de securitate minimale ale rețelei locale. Combinația MQTT +
Mosquitto este larg utilizată în proiecte IoT de cercetare \cite{naik2017mqttcomparison},
ceea ce facilitează reproductibilitatea experimentelor.

CoAP ar fi o alternativă pentru rețele cu pierderi mari sau IPv6/6LoWPAN, scenarii
irelevante pentru o rețea Wi-Fi locală. HTTP, deși simplu de implementat, ar impune un
overhead de zeci de ori mai mare față de \gls{mqtt} pentru aceeași sarcină utilă de date
senzoriale, epuizând rapid memoria stack-ului TCP de pe microcontroler.

% ======================================================================
\section{Edge computing și procesarea locală}
\label{sec:studii-edge}
% ======================================================================

Paradigma \emph{edge computing} plasează capacitatea de calcul cât mai aproape de sursa
de date, reducând latența, consumul de lățime de bandă și dependența de conectivitate la
internet. Shi et al.\ \cite{shi2016edge} definesc edge computing ca extensie a modelului
cloud la marginea rețelei și argumentează că este esențial pentru aplicații cu cerințe
stricte de latență și confidențialitate.

În contextul \gls{iot}, Premsankar et al.\ \cite{premsankar2018edgeiot} demonstrează că
procesarea locală reduce semnificativ volumul de date transmis spre cloud și permite
reacții în timp real imposibil de obținut printr-o arhitectură cu cloud central. Lin et
al.\ \cite{lin2017edge} extind analiza la scenarii de calcul mobil, arătând că un nod
edge (cum ar fi un Raspberry Pi) poate prelua sarcini de inferență \gls{ml} de complexitate
medie cu latențe sub o secundă.

Sistemul proiectat în această lucrare adoptă integral filozofia edge computing: Raspberry
Pi 3B+ acționează ca nod edge care stochează toate datele local (PostgreSQL + TimescaleDB),
execută inferența Isolation Forest în timp real și publică alerte fără a implica niciun
serviciu extern. Această arhitectură garantează funcționarea completă a sistemului chiar
și în absența totală a internetului, răspunzând direct preocupărilor de autonomie și
confidențialitate identificate în literatura de specialitate.

% ======================================================================
\section{Detecția de anomalii comportamentale}
\label{sec:studii-anomalii}
% ======================================================================

Detecția de anomalii reprezintă procesul de identificare a observațiilor care se abat
semnificativ de la comportamentul așteptat al unui sistem. Chandola et al.\
\cite{chandola2009anomaly} oferă o taxonomie exhaustivă a metodelor existente, grupate
în funcție de tipul de supraveghere și de tehnica statistică utilizată.

\subsection{Taxonomia metodelor de detecție}

\paragraph{Metode statistice.} Presupun că datele normale urmează o distribuție
probabilistică parametrică (Gaussiană, Poisson etc.). Punctele cu probabilitate mică
sunt considerate anomalii. Avantajul principal este interpretabilitatea matematică,
dar metodele parametrice eșuează când distribuția reală a datelor este necunoscută
sau multimodală — situație frecventă în datele senzoriale din locuințe.

\paragraph{Metode bazate pe distanță și densitate.} \emph{Local Outlier Factor} (LOF),
propus de Breunig et al.\ \cite{breunig2000lof}, compară densitatea locală a unui punct
cu cea a vecinilor săi. LOF este eficient pentru clustere de densități diferite, dar
costul computațional $\mathcal{O}(n^2)$ îl face nepotrivit pentru inferență în timp real
pe dispozitive edge cu resurse limitate.

\paragraph{One-Class SVM.} Scholkopf et al.\ \cite{scholkopf2001ocsvm} propun o variantă
de \gls{svm} antrenată exclusiv pe date normale, care delimitează o frontieră de decizie
în spațiul de caracteristici. Performanța depinde critic de alegerea kernelului și a
hiperparametrilor, iar costul de inferență crește cu numărul de vectori suport.

\paragraph{Metode de izolare.} Liu, Ting și Zhou \cite{liu2008isolation} introduc
\emph{Isolation Forest}, un algoritm nesupervizat care izolează punctele anomale prin
partiționări aleatoare succesive ale spațiului de caracteristici. Ideea centrală este
că un punct anomal — rar și îndepărtat — va fi separat de restul printr-un număr mic
de partiții, în timp ce un punct normal, aflat în zona densă, necesită multe partiții
pentru a fi izolat. Scorul de anomalie este invers proporțional cu lungimea medie a
căii de izolare. Liu et al.\ \cite{liu2012isolationtkdd} confirmă eficiența algoritmului
pe seturi de date de mari dimensiuni, cu complexitate $\mathcal{O}(n \log n)$ la
antrenare și $\mathcal{O}(\log n)$ la inferență.

\paragraph{Metode de învățare profundă.} Chalapathy și Chawla \cite{chalapathy2019deepanomaly}
trec în revistă utilizarea rețelelor neuronale profunde (autoencoder, LSTM, GAN) pentru
detecția de anomalii. Aceste metode ating performanțe superioare pe seturi mari de date
cu structuri complexe, dar necesită resurse computaționale semnificative (GPU, memorie
RAM de ordinul gigaocteților) și date etichetate sau volume mari de antrenare, cerințe
incompatibile cu un nod edge de tip Raspberry Pi 3B+.

Cook et al.\ \cite{cook2020anomalyiot} și Fahim și Sillitti \cite{fahim2019anomaly}
examinează specific tehnicile de detecție a anomaliilor în fluxuri \gls{iot}, subliniind
necesitatea unor metode incrementale sau cu footprint redus. Concluzia lor converge cu
alegerea Isolation Forest pentru sisteme cu constrângeri de resurse.

\subsection{Principiul izolării — ilustrare conceptuală}

\figref{fig:ml-isolation} ilustrează intuitiv diferența dintre comportamentul unui punct
normal și cel al unui punct anomal în algoritmul Isolation Forest. Un punct anomal,
aflat în afara norului dens de observații normale, este separat prin numai câteva
partiționări aleatoare. Prin contrast, un punct normal, înconjurat de alți vecini apropiați,
necesită un număr mult mai mare de partiții înainte de a fi izolat. Această asimetrie
stă la baza scorului de anomalie: cu cât calea de izolare este mai scurtă, cu atât
probabilitatea ca observația să fie o anomalie este mai mare.

\begin{figure}[H]
  \centering
  \includegraphics[width=\doublefigure]{\pics/fig-ml-isolation}
  \caption{Principiul izolării: un punct anomal este separat prin puține partiționări
    aleatoare, în timp ce un punct normal (din zona densă) necesită multe partiții.}
  \label{fig:ml-isolation}
\end{figure}

\subsection{Justificarea alegerii Isolation Forest pentru sistem}

Alegerea Isolation Forest ca algoritm de bază al motorului \gls{ml} din această lucrare
este motivată de trei proprietăți esențiale, deosebit de relevante pentru mediile edge:

\begin{enumerate}
  \item \textbf{Nesupervizat.} Nu necesită date etichetate cu anomalii cunoscute — o
    cerință imposibil de îndeplinit în practică pentru o locuință nouă, unde anomaliile
    reale sunt rare și nedefinite a priori.

  \item \textbf{Eficiență computațională.} Complexitatea $\mathcal{O}(\log n)$ la
    inferență \cite{liu2012isolationtkdd} permite scorarea în timp real a fiecărui set
    de citiri senzoriale pe un Raspberry Pi 3B+ (ARMv8, 1 GB RAM), fără a afecta
    restul serviciilor (broker MQTT, baza de date, \gls{api}).

  \item \textbf{Consum redus de memorie.} Modelul serializat (\texttt{isolation\_forest.pkl})
    ocupă câțiva zeci de kiloocteți, spre deosebire de rețelele neuronale profunde care
    pot necesita sute de megaocteți. Implementarea din scikit-learn \cite{pedregosa2011scikit}
    este matură, testată exhaustiv și bine documentată.
\end{enumerate}

Pe lângă aceste argumente tehnice, Isolation Forest s-a dovedit robust pe date
multidimensionale cu distribuții arbitrare \cite{liu2008isolation}, ceea ce este important
deoarece vectorul de caracteristici al sistemului combină valori continue (temperatură,
umiditate, gaz), valori normalizate logaritmic (luminozitate) și variabile binare/ciclice
(oră din zi, zi din săptămână, indicator noapte/weekend).

% ======================================================================
\section{Poziționarea lucrării față de stadiul actual}
\label{sec:studii-pozitionare}
% ======================================================================

Trecerea în revistă de mai sus permite identificarea clară a nișei pe care o ocupă
sistemul propus față de soluțiile existente.

\textbf{Față de platformele comerciale} (Alexa, Google Home, HomeKit), sistemul propus
elimină dependența de cloud și garantează că niciun datum personal nu părăsește rețeaua
locală. Funcționarea completă în regim offline și costul de infrastructură zero (nu există
abonamente sau servicii de hosting) constituie avantaje practice concrete pentru utilizatorul
final, în spiritul principiilor local-first formalizate de Kleppmann et al.\
\cite{kleppmann2019localfirst}.

\textbf{Față de platformele open-source generice} (Home Assistant, openHAB), sistemul
propus integrează nativ un motor de detecție a anomaliilor comportamentale antrenat pe
datele specifice locuinței, nu un sistem de reguli fix sau o integrare externă opțională.
Această abordare permite identificarea comportamentelor anormale subtile — de exemplu, un
model de temperatură neobișnuit la o anumită oră din zi — imposibil de captat prin simple
praguri hardcodate.

\textbf{Față de lucrările de cercetare pe detecția de anomalii IoT} \cite{cook2020anomalyiot,
fahim2019anomaly}, contribuția principală a prezentei lucrări constă în implementarea
completă end-to-end a unui pipeline funcțional pe hardware real de cost redus: de la
colectarea datelor senzoriale printr-un nod ESP32, prin stocarea în TimescaleDB, la
antrenarea zilnică automată a unui model Isolation Forest și inferența în timp real pe
Raspberry Pi 3B+. Lucrările citate tratează frecvent problema detecției de anomalii
în izolare, fără a aborda provocările de integrare sistem sau constrângerile edge specifice.

\textbf{Noutatea esențială} a lucrării este combinarea a trei caracteristici rareori
întâlnite simultan în literatura de specialitate: (1)~arhitectură complet \emph{local-first},
fără nicio dependență de servicii externe; (2)~detecție de anomalii comportamentale
prin \gls{ml} nesupervizat (\emph{Isolation Forest}), antrenat pe profilul real al
locuinței; (3)~implementare pe hardware de consum redus și cost accesibil (ESP32 + Raspberry
Pi), cu o interfață mobilă React Native pentru utilizatorul final.

Capitolul următor (\chapref{ch:metodologie}) detaliază arhitectura sistemului și
deciziile de proiectare care decurg din analiza prezentată mai sus.

% cap3_metodologie.tex
\chapter{Metodologie și arhitectura sistemului}
\label{ch:metodologie}

Capitolul de față descrie principiile de proiectare care au ghidat construcția sistemului,
arhitectura de ansamblu pe cinci straturi, cerințele funcționale și nefuncționale, fluxul
de date bazat pe \gls{mqtt}, modelul relațional al bazei de date, alegerile tehnologice
justificate și, în final, metodologia contribuției originale — pipeline-ul de detectare
a anomaliilor comportamentale cu algoritmul Isolation Forest.

% ─────────────────────────────────────────────────────────────────────────────
\section{Principii de proiectare}
\label{sec:principii}
% ─────────────────────────────────────────────────────────────────────────────

Arhitectura sistemului a fost concepută respectând trei principii fundamentale care
influențează fiecare decizie tehnologică.

\textbf{Local-first.} Kleppmann et al.\ definesc paradigma \emph{local-first} ca
proiectarea în care toate datele utilizatorului sunt stocate și procesate pe dispozitivul
sau infrastructura locală, conexiunea la cloud reprezentând o opțiune auxiliară și nu o
cerință de funcționare~\cite{kleppmann2019localfirst}. Sistemul implementat respectă
integral această paradigmă: senzorarea, stocarea în \gls{iot}, inferența \gls{ml} și
controlul actuatoarelor se desfășoară exclusiv în rețeaua locală (\gls{lan}), fără niciun
trafic spre servicii externe. Consecința directă este că sistemul continuă să funcționeze
normal chiar dacă conexiunea la Internet este absentă.

\textbf{Privacy by design.} Principiile \emph{privacy by design}, consolidate în
Regulamentul European \gls{gdpr}~\cite{gdpr2016}, impun ca confidențialitatea să fie
integrată în arhitectură de la etapa de proiectare, și nu adăugată ulterior. Datele de
senzori — temperaturi, prezență, niveluri de gaz — sunt date cu caracter personal în
sensul GDPR atunci când permit identificarea activității ocupanților. Prin stocarea
exclusiv locală, sistemul elimină riscul de scurgere de date prin furnizori de cloud
terți, respectând principiul minimizării transferului de date.

\textbf{Edge intelligence.} Shi et al.\ propun modelul de calcul \emph{edge computing}
ca alternativă la centralizarea în cloud, prin deplasarea puterii de calcul cât mai aproape
de sursa datelor~\cite{shi2016edge}. În sistemul de față, motorul \gls{ml} (Isolation Forest)
rulează pe Raspberry Pi 3B+ — dispozitivul gateway — și nu pe un server cloud, obținând
latență de inferență de ordinul milisecundelor și independență operațională completă față
de Internet.

% ─────────────────────────────────────────────────────────────────────────────
\section{Arhitectura de ansamblu}
\label{sec:arhitectura}
% ─────────────────────────────────────────────────────────────────────────────

Sistemul este structurat în cinci straturi ierarhice, de la hardware la interfața
utilizator. \figref{fig:arhitectura} prezintă arhitectura de ansamblu; paragrafele
de mai jos descriu fiecare strat.

\begin{figure}[H]
  \centering
  \includegraphics[width=\singlefigure]{\pics/fig-arhitectura}
  \caption{Arhitectura de ansamblu a sistemului, organizată pe cinci straturi.
           Toate componentele operează exclusiv în rețeaua locală, fără dependență de cloud.}
  \label{fig:arhitectura}
\end{figure}

\textbf{Stratul 1 — Senzori și actuatori.}
Stratul fizic cuprinde senzorii de mediu și actuatoarele conectate la nodurile ESP32.
Senzorii măsoară: temperatura și umiditatea (2× DHT11), concentrația de gaze combustibile
(MQ-2, ieșire analogică), nivelul de luminozitate (2× \gls{ldr}) și prezența umană (un
singur senzor PIR, montat pe camera exterioară). Actuatoarele includ: 3 \gls{led}-uri
de iluminat, un ventilator de 5V, un servomotor SG90 (pentru controlul draperiilor) și
un buzzer activ pentru alerte sonore locale. Camera OV2640 (montată pe ESP32-CAM)
furnizează flux video MJPEG.

\textbf{Stratul 2 — Noduri ESP32.}
Două plăci de dezvoltare gestionează interfața cu hardware-ul fizic:
\begin{itemize}
  \item \emph{Nodul interior} (ESP32 WROOM-32) execută firmware Arduino C++ și comunică
    prin \gls{gpio}, \gls{adc} (10 biți, 0–1023) și \gls{pwm} cu toți senzorii și actuatoarele
    de interior.
  \item \emph{Nodul exterior} (ESP32-CAM, AI-Thinker, cu senzor de imagine OV2640) detectează
    mișcarea prin PIR pe GPIO13 și oferă server HTTP pe portul 80 pentru streaming video
    (\texttt{GET /stream}) și captură JPEG (\texttt{GET /capture}).
\end{itemize}
Ambele noduri comunică cu gateway-ul prin \gls{mqtt} peste Wi-Fi; camera este descoperită
de aplicația mobilă prin \gls{mdns} (\texttt{smarthome-cam.local}).

\textbf{Stratul 3 — Gateway Raspberry Pi.}
Raspberry Pi 3B+ îndeplinește rolul central al sistemului. Găzduiește:
\begin{itemize}
  \item brokerul Mosquitto (\gls{mqtt} 1883, autentificat), care primește și rutează
    toate mesajele publicate de noduri;
  \item baza de date PostgreSQL cu extensia TimescaleDB, pentru stocarea eficientă
    a seriilor temporale de citiri;
  \item motorul \gls{ml} Python (Isolation Forest), care antrenează și rulează inferența
    de detectare a anomaliilor direct pe dispozitivul edge.
\end{itemize}

\textbf{Stratul 4 — \gls{api} NestJS.}
Serviciul backend, scris în TypeScript cu framework-ul NestJS, expune un \gls{api} \gls{rest}
și un server WebSocket (Socket.IO) pe portul 3000. Comunicarea cu brokerul \gls{mqtt} se
realizează prin modulul \texttt{mqtt.service.ts}, care abonează NestJS la topicurile
de senzori și publică comenzi spre noduri. Autentificarea este realizată cu \gls{jwt}
și parole \emph{hash}-uite cu bcrypt.

\textbf{Stratul 5 — Aplicație mobilă.}
Interfața utilizatorului este o aplicație React Native (Expo SDK 54, TypeScript) ce
rulează pe Android și iOS. Primește date în timp real prin Socket.IO, afișează
valorile senzorilor, starea alertelor și permite controlul actuatoarelor. Descoperirea
automată a gateway-ului se face prin \gls{mdns} (\texttt{smarthome.local}).

% ─────────────────────────────────────────────────────────────────────────────
\section{Cerințe funcționale și nefuncționale}
\label{sec:cerinte}
% ─────────────────────────────────────────────────────────────────────────────

\subsection{Cerințe funcționale}

Cerințele funcționale au fost identificate pornind de la scenariile de utilizare reale
ale unui locatar. \figref{fig:usecase} prezintă diagrama use-case a sistemului.

\begin{figure}[H]
  \centering
  \includegraphics[width=\singlefigure]{\pics/fig-usecase}
  \caption{Diagrama use-case a sistemului smart home. Actorul principal este utilizatorul,
           iar motorul ML participă ca actor secundar (sistem) prin generarea automată
           de alerte de anomalie.}
  \label{fig:usecase}
\end{figure}

Principalele cerințe funcționale sunt:

\begin{enumerate}
  \item \textbf{Monitorizare senzori în timp real} — sistemul colectează și afișează
    temperaturi (2 camere), umiditate (2 camere), nivel de gaz, luminozitate (2 senzori)
    și prezență, cu actualizare la fiecare 5 secunde.
  \item \textbf{Control actuatori de la distanță} — utilizatorul poate acționa de pe
    telefon: LED-urile (pornit/oprit/toggle individual), ventilatorul, servomotorul
    draperiilor (0°/90°/180°), buzzer-ul și comanda \texttt{all\_off}.
  \item \textbf{Sistem de alerte} — alertele de tip \texttt{GAS\_DETECTED},
    \texttt{MOTION\_NIGHT}, \texttt{ALERT\_HEAT}, \texttt{ALERT\_COLD} și
    \texttt{ML\_ANOMALY} sunt stocate în baza de date și notificate în timp real
    prin WebSocket.
  \item \textbf{Vizualizare istoric} — aplicația permite consultarea istoricului
    grafic al citirilor pe intervale configurabile.
  \item \textbf{Flux video exterior} — aplicația afișează stream MJPEG live și
    captură snapshot de la camera exterioară.
  \item \textbf{Detectare anomalii comportamentale} — motorul \gls{ml} local
    generează alerte de tip \texttt{ML\_ANOMALY} pentru comportamente neobișnuite
    față de pattern-ul zilnic al locatarului.
  \item \textbf{Configurare Wi-Fi (provisioning)} — nodul poate primi noile
    credențiale Wi-Fi prin comanda \gls{mqtt} \texttt{wifi\_update} sau, ca
    fallback, prin hotspot-ul \texttt{SmartHome-Setup}.
\end{enumerate}

\subsection{Cerințe nefuncționale}

\begin{itemize}
  \item \textbf{Disponibilitate locală} — sistemul funcționează integral fără
    conexiune la Internet; nu există dependențe de servicii cloud externe.
  \item \textbf{Latență redusă} — de la publicarea datelor de senzori până la
    afișarea în aplicație, latența end-to-end vizată este sub 2 secunde.
  \item \textbf{Confidențialitate} — nicio dată de senzori sau video nu
    părăsește rețeaua locală.
  \item \textbf{Extensibilitate} — arhitectura permite adăugarea de noi noduri
    ESP32 fără modificări în backend, prin convenția de topicuri
    \texttt{smarthome\slb{} \{nodeId\}\slb{}sensors}.
  \item \textbf{Reziliență} — nodurile ESP32 implementează reconectare automată
    la Wi-Fi și \gls{mqtt}; LWT asigură notificarea automată a deconectărilor.
\end{itemize}

% ─────────────────────────────────────────────────────────────────────────────
\section{Fluxul de date și protocolul MQTT}
\label{sec:flux-date}
% ─────────────────────────────────────────────────────────────────────────────

\gls{mqtt} este protocolul de transport ales pentru comunicarea dintre noduri
și gateway. Standardul \gls{oasis} MQTT 3.1.1~\cite{oasis2019mqtt} definește un
protocol publish/subscribe binar, cu overhead minim, potrivit pentru dispozitive
cu resurse limitate~\cite{naik2017mqttcomparison}. Brokerul utilizat este
Mosquitto~\cite{light2017mosquitto}, rulând cu autentificare bazată pe parolă
pe portul 1883.

\figref{fig:flux-date} prezintă fluxul complet de date din sistem.

\begin{figure}[H]
  \centering
  \includegraphics[width=\singlefigure]{\pics/fig-flux-date}
  \caption{Fluxul de date în sistem. Datele de senzori circulă de la ESP32 prin
           Mosquitto spre NestJS, care le persistă în TimescaleDB, le transmite
           aplicației prin WebSocket și le evaluează prin motorul ML.
           Comenzile urmează traseul invers (linie punctată).}
  \label{fig:flux-date}
\end{figure}

\subsection{Topicuri MQTT și contractul de date}

\tabref{tab:mqtt-topics} detaliază topicurile \gls{mqtt} utilizate, direcția
fluxului, nivelul QoS și rolul fiecăruia.

\begin{table}[H]
  \centering
  \caption{Topicurile MQTT ale sistemului smart home.}
  \label{tab:mqtt-topics}
  \small
  \begin{tabular}{@{}p{4.4cm} p{2.4cm} p{1.5cm} p{4.6cm}@{}}
    \toprule
    \textbf{Topic} & \textbf{Direcție} & \textbf{QoS} & \textbf{Descriere} \\
    \midrule
    \texttt{smarthome\slb{} \{nodeId\}\slb{}sensors}
      & ESP32 → broker
      & 0
      & Citiri periodice (la 5 s) în format JSON; payload conține
        temperaturi, umiditate, gaz, luminozitate. \\
    \addlinespace
    \texttt{smarthome\slb{} \{nodeId\}\slb{}status}
      & ESP32 → broker
      & 1, retained
      & Starea nodului: \texttt{online} la conectare,
        \texttt{offline} prin Last Will (LWT). \\
    \addlinespace
    \texttt{smarthome\slb{} \{nodeId\}\slb{}commands}
      & broker → ESP32
      & 1
      & Comenzi pentru actuatori: \texttt{led\_on}, \texttt{led\_off},
        \texttt{led\_toggle}, \texttt{fan\_on}, \texttt{fan\_off},
        \texttt{servo\_move}, \texttt{buzzer\_beep}, \texttt{all\_off},
        \texttt{wifi\_update}. \\
    \addlinespace
    \texttt{smarthome/alerts}
      & ESP32 / ML → broker
      & 1
      & Alerte de urgență: gaz detectat (de pe nod) sau anomalie ML
        (de pe gateway). \\
    \bottomrule
  \end{tabular}
\end{table}

Payloadul publicat de nodul interior pe topicul de senzori este un obiect JSON
cu câmpurile prezentate în fragmentul de mai jos (contract comun între firmware,
backend și aplicația mobilă):

\begin{lstlisting}[language=TypeScript, caption={Payload MQTT publicat de nodul interior la fiecare 5 secunde.}]
{
  "nodeId":       "esp32_node_a",
  "location":     "interior",
  "temperature":  22.5,   // DHT11 #1, living (°C)
  "humidity":     55,     // DHT11 #1, living (%RH)
  "temperature2": 21.8,   // DHT11 #2, dormitor (°C)
  "humidity2":    58,     // DHT11 #2, dormitor (%RH)
  "gasLevel":     245,    // MQ-2 ADC 0-1023
  "gasAlert":     false,  // true dacă gasLevel > 350
  "light1":       410,    // LDR #1 ADC 0-1023
  "light2":       120,    // LDR #2 ADC 0-1023
  "lightLux":     410     // alias light1 (compatibilitate)
}
\end{lstlisting}

\subsection{Last Will and Testament și QoS}

La conectare, nodul configurează mesajul \emph{Last Will and Testament} (LWT) pe topicul
\texttt{smarthome\slb{} \{nodeId\}\slb{}status} cu payload \texttt{\{"status":"offline"\}} și flag
\emph{retained}. Dacă conexiunea TCP se întrerupe fără deconectare explicită, brokerul
publică automat mesajul LWT, permițând NestJS să marcheze nodul ca offline în baza
de date și să notifice aplicația mobilă prin WebSocket.

Datele de senzori se publică cu QoS~0 (fire-and-forget), acceptabil pentru telemetrie
continuă cu cadență mare (pierderea unui pachet din 720 de zi este neglijabilă).
Comenzile și alertele folosesc QoS~1 (cel puțin o dată), asigurând livrarea
confirmată.

\subsection{Descoperire prin mDNS și provisioning Wi-Fi}

Descoperirea automată a broker-ului și a camerei utilizează \gls{mdns} (RFC 6762~\cite{rfc6762mdns}):
ESP32 rezolvă \texttt{smarthome.local} la adresa IP curentă a Raspberry Pi, fără
configurare manuală. Camera exterioară este disponibilă la \texttt{smarthome-cam.local}.

Provisioningul Wi-Fi urmează o schemă cu fallback: credențialele curente (SSID, parolă,
adresa broker) sunt persistate în memoria NVS (\emph{Non-Volatile Storage}) a ESP32.
La pornire, nodul încearcă conectarea la rețeaua salvată; în caz de eșec, revine automat
la hotspot-ul de configurare \texttt{SmartHome-Setup} (IP fix \texttt{192.168.4.1}).
Credențialele pot fi actualizate în orice moment prin comanda \gls{mqtt}
\texttt{wifi\_update \{ssid, password, mqtt\_host\}}, după care nodul repornește
și se reconectează la noua rețea.

% ─────────────────────────────────────────────────────────────────────────────
\section{Modelul de date}
\label{sec:model-date}
% ─────────────────────────────────────────────────────────────────────────────

Baza de date rulează pe PostgreSQL~\cite{postgresql2024} cu extensia
TimescaleDB~\cite{timescaledb2017}, specializată în stocarea și interogarea eficientă
a seriilor temporale~\cite{jensen2017timeseries}. \figref{fig:er} prezintă diagrama
entitate-relație.

\begin{figure}[H]
  \centering
  \includegraphics[width=\singlefigure]{\pics/fig-er}
  \caption{Diagrama entitate-relație a bazei de date. Tabelul \texttt{sensor\_readings}
           este hypertable TimescaleDB particionat automat pe coloana \texttt{time}.}
  \label{fig:er}
\end{figure}

\textbf{Tabelul \texttt{sensor\_readings} (hypertable).}
Este tabelul central al sistemului, convertit în hypertable TimescaleDB cu partiționare
automată pe coloana \texttt{time}. Stochează fiecare citire publicată de un nod ESP32:

\begin{itemize}
  \item \texttt{time} — \texttt{TIMESTAMPTZ}, coloana de partiționare;
  \item \texttt{node\_id}, \texttt{location} — identificatorul și locația nodului;
  \item \texttt{temperature}, \texttt{humidity} — citirile DHT11~\#1 (living);
  \item \texttt{temperature2}, \texttt{humidity2} — citirile DHT11~\#2 (dormitor),
    \texttt{NULLABLE} (câmpuri absente în citiri mai vechi sau în caz de eroare senzor);
  \item \texttt{gas\_level}, \texttt{gas\_alert} — valoarea \gls{adc} MQ-2 și indicatorul
    de alertă;
  \item \texttt{light1}, \texttt{light2}, \texttt{light\_lux} — valorile celor două
    \gls{ldr}-uri și alias-ul de compatibilitate.
\end{itemize}

TimescaleDB comprimă automat chunk-urile mai vechi de 7 zile și permite interogări
de agregare temporală cu funcții native (\texttt{time\_bucket}, \texttt{first},
\texttt{last}) considerabil mai rapide decât PostgreSQL simplu pe seturi mari
de date~\cite{timescaledb2017}.

\textbf{Tabelul \texttt{alerts}.}
Stochează alertele generate, cu câmpurile: \texttt{id}, \texttt{time}, \texttt{node\_id},
\texttt{alert\_type} (unul dintre \texttt{GAS\_DETECTED}, \texttt{MOTION\_NIGHT},
\texttt{ALERT\_HEAT}, \texttt{ALERT\_COLD}, \texttt{ML\_ANOMALY}),
\texttt{severity}, \texttt{location}, \texttt{details} (JSON cu context),
\texttt{acknowledged} (boolean).

\textbf{Tabelul \texttt{nodes}.}
Registrul nodurilor active: \texttt{id}, \texttt{node\_id}, \texttt{location},
\texttt{last\_seen}, \texttt{online}, \texttt{node\_type}.

\textbf{Tabelul \texttt{automation\_rules}.}
Permite definirea de reguli de automatizare configurabile: senzor, operator,
prag (\texttt{threshold}), interval orar (\texttt{time\_from}, \texttt{time\_to}),
tip de acțiune (\texttt{action\_type}), stare activă.

\textbf{Tabelele \texttt{users} și \texttt{ml\_models}.}
\texttt{users} păstrează credențialele autentificate (username, \texttt{password\_hash}
produs prin bcrypt). \texttt{ml\_models} înregistrează metadatele fiecărui model antrenat:
versiune (formatul \texttt{vAAAAMMZZ\_HHMM}), numărul de rânduri de antrenare,
valoarea \texttt{contamination}, lista feature-urilor și câmpul \texttt{is\_active}
care identifică modelul curent utilizat în inferență.

% ─────────────────────────────────────────────────────────────────────────────
\section{Alegeri tehnologice justificate}
\label{sec:tehnologii}
% ─────────────────────────────────────────────────────────────────────────────

\tabref{tab:tehnologii} sintetizează principalele alegeri tehnologice și justificarea
lor în raport cu cerințele sistemului.

\begin{table}[H]
  \centering
  \caption{Alegeri tehnologice — cerință, soluție adoptată și justificare.}
  \label{tab:tehnologii}
  \begin{tabular}{@{}p{3.5cm} p{3.2cm} p{6.5cm}@{}}
    \toprule
    \textbf{Cerință} & \textbf{Tehnologie} & \textbf{Justificare} \\
    \midrule
    Microcontroler \gls{iot} cu Wi-Fi
      & ESP32 WROOM-32~\cite{espressif2023esp32}
      & Dual-core 240~MHz, 520~KB SRAM, Wi-Fi 802.11b/g/n integrat, suport Arduino;
        cost redus (\textasciitilde5~USD); \gls{adc} 12-biți pe mai mulți pini. \\
    \addlinespace
    Transport mesaje pub/sub
      & MQTT 3.1.1 / Mosquitto~\cite{oasis2019mqtt,light2017mosquitto}
      & Protocol binar cu overhead minim (\textasciitilde2~B header); QoS 0–2;
        LWT pentru detectarea deconectărilor; potrivit pentru rețele locale cu
        lățime de bandă redusă. \\
    \addlinespace
    Stocare serii temporale
      & PostgreSQL + TimescaleDB~\cite{timescaledb2017,postgresql2024}
      & Partiționare automată (hypertable) pe timp; compresie nativă; funcții de
        agregare temporală optimizate; interoperabilitate SQL deplină cu TypeORM. \\
    \addlinespace
    Backend \gls{api}
      & NestJS + TypeScript~\cite{nestjs2024}
      & Arhitectură modulară cu injecție de dependențe; integrare nativă WebSocket
        (Socket.IO); TypeORM pentru acces la baza de date; suport \gls{jwt} din cutie. \\
    \addlinespace
    Aplicație mobilă
      & React Native / Expo~\cite{reactnative2024,expo2024}
      & Cod TypeScript partajat între iOS și Android; Expo SDK 54 simplifică
        build-ul; ecosistem bogat (React Query, Zustand, Socket.IO client). \\
    \addlinespace
    Detectare anomalii
      & Isolation Forest~\cite{liu2008isolation,liu2012isolationtkdd}
      & Algoritm nesupervizat, fără etichete; complexitate de antrenare
        $O(n \log n)$; consum de memorie redus (\textasciitilde4~MB model pkl);
        implementare matură în scikit-learn~\cite{pedregosa2011scikit};
        rulabil pe Raspberry Pi 3B+ fără accelerator hardware. \\
    \addlinespace
    Autentificare
      & JWT + bcrypt~\cite{jones2015jwt,provos2003bcrypt}
      & Tokenuri fără stare, ușor de validat la fiecare request \gls{rest};
        bcrypt cu factor de cost adaptabil pentru hash-uirea parolelor. \\
    \bottomrule
  \end{tabular}
\end{table}

% ─────────────────────────────────────────────────────────────────────────────
\section{Pipeline-ul de detectare a anomaliilor — contribuția originală}
\label{sec:ml-pipeline}
% ─────────────────────────────────────────────────────────────────────────────

Contribuția originală centrală a lucrării este un pipeline de detectare a anomaliilor
comportamentale bazat pe algoritmul Isolation Forest, implementat în Python și rulând
local pe Raspberry Pi 3B+. Spre deosebire de sistemele care transmit datele la cloud
pentru analiză, abordarea propusă menține toate datele și calculele în rețeaua locală,
respectând principiile de confidențialitate enunțate în \secref{sec:principii}.

\figref{fig:ml-pipeline} prezintă structura completă a pipeline-ului.

\begin{figure}[H]
  \centering
  \includegraphics[width=\singlefigure]{\pics/fig-ml-pipeline}
  \caption{Pipeline-ul complet al motorului ML: de la colectarea datelor din
           TimescaleDB, prin ingineria trăsăturilor și normalizare, până la
           antrenarea periodică a modelului Isolation Forest și inferența în timp real.
           Bucla de reantrenare zilnică și mecanismul de cold-start sunt evidențiate.}
  \label{fig:ml-pipeline}
\end{figure}

\subsection{Algoritmul Isolation Forest}

Isolation Forest a fost propus de Liu, Ting și Zhou~\cite{liu2008isolation} și
extins ulterior~\cite{liu2012isolationtkdd} ca algoritm de detectare a anomaliilor
bazat pe proprietatea că punctele anomale sunt mai ușor de „izolat" prin partiționări
aleatoare ale spațiului de trăsături față de punctele normale. Formal, un arbore
de izolare (\emph{iTree}) partiționează recursiv un subsamplu de date prin selecție
aleatoare a unei trăsături și a unui prag de separare. Lungimea medie a căii de la
rădăcină până la un punct $\mathbf{x}$ în toți arborii constituie scorul de anomalie.

Scorul de anomalie $s(\mathbf{x}, n)$ este definit ca:

\begin{equation}
  s(\mathbf{x}, n) = 2^{-\dfrac{E[h(\mathbf{x})]}{c(n)}},
  \label{eq:if-score}
\end{equation}

unde $h(\mathbf{x})$ este lungimea căii parcurse de $\mathbf{x}$ în arbore, $E[\cdot]$
este media peste toți arborii din ansamblu, $n$ este dimensiunea subsamplului de antrenare
și $c(n) = 2H(n-1) - \tfrac{2(n-1)}{n}$ este lungimea medie a căii pentru $n$ noduri
($H$ — număr armonic). Un scor $s \to 1$ indică o anomalie pronunțată (cale scurtă),
iar $s \to 0$ indică un punct normal. Implementarea din scikit-learn~\cite{pedregosa2011scikit}
returnează \texttt{score\_samples(x)} ca valori negative, cu semnatica: cu cât scorul
este mai negativ, cu atât punctul este mai anormal.

\subsection{Ingineria trăsăturilor (\texttt{features.py})}

Modulul \texttt{features.py} construiește 12 trăsături numerice din citirile brute ale
senzorilor, în ordinea fixă din \tabref{tab:features}. Ordinea este identică în
antrenare și în inferență — orice modificare necesită reantrenare completă a modelului.

\begin{table}[H]
  \centering
  \caption{Cele 12 trăsături ale modelului Isolation Forest, în ordinea fixă din \texttt{FEATURE\_NAMES}.}
  \label{tab:features}
  \begin{tabular}{@{}r l p{7.5cm}@{}}
    \toprule
    \textbf{\#} & \textbf{Trăsătură} & \textbf{Descriere și sursă} \\
    \midrule
    0  & \texttt{temperature}      & Temperatura DHT11~\#1 (living), °C. \\
    1  & \texttt{humidity}         & Umiditatea relativă DHT11~\#1 (\%RH). \\
    2  & \texttt{gas\_normalized}  & Valoarea ADC MQ-2 împărțită la 4095
                                     (normalizare în $[0,1]$). \\
    3  & \texttt{lux\_normalized}  & $\log_{1p}(\texttt{light1}) / 7{,}0$ —
                                     luminozitate LDR~\#1 normalizată logaritmic
                                     pe scala ADC 0–1023
                                     ($\log_{1p}(1023) \approx 6{,}93$). \\
    4  & \texttt{motion}           & Prezența umană detectorului PIR (0/1). \\
    5  & \texttt{hour\_of\_day}    & Ora din zi (0–23), captează ritmul circadian. \\
    6  & \texttt{day\_of\_week}    & Ziua săptămânii (0 = luni, 6 = duminică). \\
    7  & \texttt{is\_weekend}      & 1 dacă ziua este sâmbătă sau duminică. \\
    8  & \texttt{is\_night}        & 1 dacă ora $\geq 23$ sau ora $< 6$. \\
    9  & \texttt{temp\_delta\_1h}  & Diferența de temperatură față de citirile
                                     cu 2 eșantioane în urmă (aproximativ 10~s). \\
    10 & \texttt{motion\_count\_10}& Suma evenimentelor de mișcare în ultimele
                                     10 citiri (fereastră glisantă). \\
    11 & \texttt{temp\_avg}        & Media celor doi senzori DHT11; dacă DHT11~\#2
                                     lipsește, se folosește valoarea DHT11~\#1. \\
    \bottomrule
  \end{tabular}
\end{table}

\subsection{Normalizare}

Înainte de antrenare, trăsăturile sunt normalizate cu \texttt{StandardScaler} din
scikit-learn, care transformă fiecare trăsătură la medie 0 și deviație standard 1:
\begin{equation}
  z_i = \frac{x_i - \mu_i}{\sigma_i}.
  \label{eq:standardscaler}
\end{equation}
Parametrii $\mu_i$ și $\sigma_i$ sunt calculați pe datele de antrenare și salvați în
\texttt{scaler.pkl} alături de model. La inferență, aceeași transformare este aplicată
fiecărei noi citiri înainte de evaluarea scorului.

\subsection{Antrenarea modelului (\texttt{train.py})}

Modulul \texttt{train.py} implementează ciclul de antrenare:

\begin{enumerate}
  \item Citește din tabela \texttt{sensor\_readings} toate înregistrările din
    ultimele 30 de zile cu \texttt{gas\_alert = FALSE} (citeșterile din perioadele
    de alertă sunt excluse pentru a nu contamina distribuția de normalitate).
  \item Verifică pragul minim de 500 de rânduri; dacă datele sunt insuficiente,
    antrenarea este întreruptă cu mesaj informativ.
  \item Apelează \texttt{build\_features(df)} și normalizează cu \texttt{StandardScaler}.
  \item Antrenează un \texttt{IsolationForest} cu parametrii:
    \texttt{n\_estimators=100}, \texttt{contamination=0.05}, \texttt{random\_state=42},
    \texttt{n\_jobs=-1} (toate nucleele disponibile).
  \item Salvează \texttt{isolation\_forest.pkl} și \texttt{scaler.pkl} în directorul
    \texttt{models/}.
  \item Înregistrează metadatele modelului nou în tabela \texttt{ml\_models} și
    marchează toate versiunile anterioare ca inactive.
\end{enumerate}

Parametrul \texttt{contamination=0.05} afirmă că modelul estimează că 5\% din datele
de antrenare reprezintă comportamente atipice. Această valoare este conservatoare —
subestimarea duce la mai puține fals-negative în faza de deploy. Reantrenarea este
planificată zilnic la ora 03:00 prin \texttt{scheduler.py}, rulat ca serviciu systemd
\texttt{smarthome-ml.service} pe Raspberry Pi.

\subsection{Inferența în timp real (\texttt{infer.py})}

Clasa \texttt{AnomalyDetector} din \texttt{infer.py} este instanțiată la pornirea
serviciului NestJS și menținută în memorie. La fiecare citire nouă de senzori:

\begin{enumerate}
  \item Se construiește un DataFrame cu o singură linie din payload-ul JSON primit.
  \item Se aplică \texttt{build\_features} și \texttt{scaler.transform}.
  \item Se apelează \texttt{model.score\_samples(X\_scaled)}, care returnează un scalar
    real; notăm scorul ca $s \in \mathbb{R}$.
  \item Dacă $s < \tau$ (unde $\tau = -0{,}1$ este pragul \texttt{ANOMALY\_THRESHOLD}),
    citirea este clasificată drept anomalie și se creează o alertă de tip
    \texttt{ML\_ANOMALY} în baza de date, notificată imediat aplicației prin WebSocket.
\end{enumerate}

Metoda returnează un obiect cu câmpurile \texttt{is\_anomaly} (boolean), \texttt{score}
(valoare numerică), \texttt{threshold} și \texttt{model\_active}. Dacă fișierele
\texttt{.pkl} nu există, câmpul \texttt{model\_active} este \texttt{false} și
nicio alertă ML nu este generată.

\subsection{Mecanismul de cold-start}

În primele aproximativ 14 zile de funcționare a sistemului, nu există suficiente date
pentru a antrena modelul (minimul de 500 de rânduri echivalând cu circa 42 de minute
de date, dar o fereastră de 30 de zile este necesară pentru a surprinde variabilitatea
comportamentală zilnică și săptămânală). În această perioadă de cold-start, motorul
ML returnează \texttt{model\_active: false} și sistemul se bazează pe reguli de prag
hardcodate pentru generarea alertelor:

\begin{itemize}
  \item \texttt{gas\_level > 300} → alertă \texttt{GAS\_DETECTED};
  \item \texttt{temperature > 35°C} → alertă \texttt{ALERT\_HEAT};
  \item \texttt{temperature < 5°C} → alertă \texttt{ALERT\_COLD};
  \item mișcare detectată în intervalul 23:00–06:00 → alertă \texttt{MOTION\_NIGHT}.
\end{itemize}

Această strategie asigură că sistemul oferă protecție de bază chiar înainte ca modelul
\gls{ml} să fie disponibil.

% ─────────────────────────────────────────────────────────────────────────────
\section{Considerații de securitate}
\label{sec:securitate}
% ─────────────────────────────────────────────────────────────────────────────

Modelul de securitate al sistemului este construit pe trei niveluri complementare.

\textbf{Autentificare și autorizare.}
Accesul la \gls{api}-ul \gls{rest} și la conexiunile WebSocket este protejat prin
\gls{jwt}~\cite{jones2015jwt}. La autentificarea reușită, serverul emite un token semnat
cu o cheie secretă; toate request-urile ulterioare trebuie să includă tokenul în header-ul
\texttt{Authorization: Bearer <token>}. Parolele utilizatorilor sunt stocate exclusiv
ca hash-uri bcrypt~\cite{provos2003bcrypt} cu factor de cost adaptabil (implicit 10
runde), eliminând riscul expunerii parolelor în clar chiar în cazul unui acces
neautorizat la baza de date.

\textbf{Autentificare broker MQTT.}
Brokerul Mosquitto este configurat cu autentificare prin user/parolă (utilizatorul
\texttt{smarthome}, credențiale stocate în fișierul Mosquitto \texttt{passwd}),
permițând conectarea doar nodurilor și serviciilor cunoscute. Comunicarea se
desfășoară în rețeaua locală și nu este expusă direct la Internet.

\textbf{Izolare în rețeaua locală.}
Nicio componentă a sistemului nu expune porturi pe Internet. Gateway-ul Raspberry Pi,
nodurile ESP32 și aplicația mobilă comunică exclusiv în \gls{lan}. Această izolare
reprezintă o măsură de securitate perimetrală complementară celor aplicative,
conformă cu principiul \emph{privacy by design} din \gls{gdpr}~\cite{gdpr2016}.
Datele cu caracter personal (activitate comportamentală, prezență umană) nu
părăsesc niciodată rețeaua locală, eliminând riscurile asociate transferului
transfrontalier de date sau breșelor de securitate la nivel de furnizor cloud.

\textbf{Minimizarea suprafeței de atac.}
Firmware-ul ESP32 nu expune nicio interfață HTTP sau Telnet în mod normal de funcționare;
singurul vector de comunicare externă este conexiunea \gls{mqtt} autentificată.
ESP32-CAM expune server HTTP pe portul 80 exclusiv în rețeaua locală, accesul nefiind
publicat prin NAT sau port-forwarding.

% cap4_implementare.tex
\chapter{Implementarea prototipului}
\label{ch:implementare}

Capitolul de față descrie concret modul în care componentele prezentate în
\chapref{ch:metodologie} au fost transpuse în cod funcțional și asamblate
într-un prototip operațional.
Sunt tratate, în ordine, firmware-ul nodului interior bazat pe
\gls{esp32}, firmware-ul camerei exterioare \gls{esp32}-CAM, configurarea
gateway-ului Raspberry Pi, structura \gls{api}-ului NestJS, motorul de
detecție a anomaliilor cu \gls{if} și, în final, aplicația mobilă.

% ──────────────────────────────────────────────────────────────────
\section{Firmware-ul nodului interior (ESP32 WROOM-32)}
\label{sec:firmware-nod}

\subsection{Pinout și configurare hardware}
\label{sec:pinout}

Nodul interior este construit în jurul modulului ESP32 WROOM-32.
Toți pinii \gls{adc} folosiți pentru senzori (MQ-2, \gls{ldr}\,\#1 și
\gls{ldr}\,\#2) aparțin grupului ADC1 (GPIO34, GPIO35, GPIO36), ceea ce
permite citiri analogice corecte chiar și cu \gls{wifi} activ — spre
deosebire de pinii ADC2 care devin inutilizabili în prezența
radio-ului~\cite{espressif2023esp32}.
Modulul BH1750 (citire de lux prin I\textsuperscript{2}C), luat în
considerare inițial, a fost eliminat în favoarea celor două \gls{ldr}
analogice, simplificând schema de cablare.

\tabref{tab:pinout} centralizează cei unsprezece pini activi ai
nodului, împărțiți în senzori și actuatori.

\begin{table}[H]
  \centering
  \caption{Pinout complet al nodului interior (ESP32 WROOM-32).}
  \label{tab:pinout}
  \begin{tabular}{@{} l l p{6.2cm} @{}}
    \toprule
    \textbf{Componentă} & \textbf{Pin GPIO} & \textbf{Observații} \\
    \midrule
    \gls{dht}11 \#1 (living)   & GPIO4  & Digital, pull-up extern 10\,k\textOmega{} la 3{,}3\,V \\
    \gls{dht}11 \#2 (dormitor) & GPIO23 & Digital, pull-up extern 10\,k\textOmega{} la 3{,}3\,V \\
    MQ-2 (gaz, A0)             & GPIO34 & ADC1, intrare; divizor de tensiune 0–3{,}3\,V \\
    \gls{ldr} \#1 (lumină)     & GPIO35 & ADC1, intrare; divizor cu 10\,k\textOmega{} \\
    \gls{ldr} \#2 (lumină)     & GPIO36 & ADC1 (VP), intrare; divizor cu 10\,k\textOmega{} \\
    Buzzer activ               & GPIO5  & Ieșire digitală, oscilator intern \\
    Servo SG90 (draperii)       & GPIO18 & \gls{pwm} 50\,Hz; alimentat din 5\,V extern \\
    \gls{led} \#0 (living)     & GPIO25 & Ieșire digitală + rezistor serie 220\,\textOmega{} \\
    \gls{led} \#1 (dormitor)   & GPIO26 & Ieșire digitală + rezistor serie 220\,\textOmega{} \\
    \gls{led} \#2 (hol)        & GPIO27 & Ieșire digitală + rezistor serie 220\,\textOmega{} \\
    Ventilator 5\,V            & GPIO33 & Comandat prin tranzistor NPN (NU direct din GPIO) \\
    \bottomrule
  \end{tabular}
\end{table}

\figref{fig:pinout} ilustrează schema de conexiuni a nodului, cu marcarea
alimentărilor separate: 3{,}3\,V pentru senzori (\gls{dht}11, \gls{ldr})
și 5\,V extern pentru servo și ventilator.

\begin{figure}[H]
  \centering
  \includegraphics[width=\singlefigure]{\pics/fig-pinout}
  \caption{Schema de conexiuni a nodului interior ESP32 WROOM-32
           (pinout conform \tabref{tab:pinout}).}
  \label{fig:pinout}
\end{figure}

\subsection{Citirea senzorilor}
\label{sec:citire-senzori}

La inițializare, rezoluția \gls{adc} este setată explicit la 10 biți
(\texttt{analogReadResolution(10)}), oferind valori în intervalul 0–1023.
Ambii senzori \gls{dht}11 sunt citiți cu biblioteca \emph{Adafruit DHT},
iar erorile temporare (valori \texttt{NaN}) sunt tratate tolerant:
dacă cel puțin un senzor raportează date valide, publicarea continuă cu
câmpurile disponibile; numai dacă ambii eșuează simultan este omisă
întreaga publicare~\cite{aosong2019dht11}.

Senzorul de gaz MQ-2 returnează o valoare analogică brută (0–1023).
Când valoarea depășește pragul \texttt{GAS\_THRESHOLD = 350}, firmware-ul
detectează \emph{frontul crescător} (tranziția \texttt{false → true}),
publică o alertă pe topicul \texttt{smarthome/alerts} și activează
buzzerul local cu trei bipuri — prevenind astfel spamul la broker în caz
de depășire prelungită~\cite{hanwei2020mq2}.
\listingref{lst:publish-sensors} prezintă nucleul funcției de colectare
și publicare.

\begin{lstlisting}[language=C++,
  caption={Fragment din \texttt{publishSensors()} — citire DHT11, MQ-2,
           LDR și construire payload JSON (esp32\_node.ino).},
  label=lst:publish-sensors]
// Citire senzori ADC — 10 biti (0-1023)
int gas    = analogRead(MQ2_PIN);   // GPIO34
int light1 = analogRead(LDR1_PIN);  // GPIO35
int light2 = analogRead(LDR2_PIN);  // GPIO36

// Alerta gaz — doar pe frontul crescator
bool gasAlert = (gas > GAS_THRESHOLD);
if (gasAlert && !gasAlertActive) {
  publishGasAlert(gas);
  beep(3);         // semnal sonor local imediat
}
gasAlertActive = gasAlert;

// Payload JSON (contract cu backend si mobile)
StaticJsonDocument<384> doc;
doc["nodeId"]    = NODE_ID;          // "esp32_node_a"
doc["location"]  = LOCATION;        // "interior"
if (dht1ok) { doc["temperature"] = temp1; doc["humidity"] = humid1; }
if (dht2ok) { doc["temperature2"]= temp2; doc["humidity2"]= humid2; }
doc["gasLevel"]  = gas;
doc["gasAlert"]  = gasAlert;
doc["light1"]    = light1;
doc["light2"]    = light2;
doc["lightLux"]  = light1;          // alias backward-compat
\end{lstlisting}

\subsection{Actuatori}
\label{sec:actuatori}

Ventilatorul (5\,V) nu poate fi alimentat direct din \gls{gpio} —
curentul maxim al unui pin ESP32 este de 12\,mA — astfel că este comandat
printr-un tranzistor NPN pe GPIO33: nivel HIGH pe pin conduce tranzistorul
și închide circuitul motorului.
Servo-ul SG90, folosit pentru controlul draperiilor, primește semnal
\gls{pwm} la 50\,Hz prin biblioteca \emph{ESP32Servo}; funcția
\texttt{servoDraperii.write(angle)} acceptă unghiuri între 0° și 180°
(corespunzând draperiilor complet deschise, respectiv complet închise).

Ventilatorul poate fi acționat în două moduri, comutabile din aplicație:
\emph{manual} (comenzile \texttt{fan\_on}/\texttt{fan\_off}) și
\emph{automat} (comanda \texttt{fan\_auto}). În modul automat, firmware-ul
pornește ventilatorul când temperatura medie a celor doi senzori \gls{dht}11
depășește un \textbf{prag configurabil de la distanță} (comanda
\texttt{fan\_threshold}, implicit 28\,°C) și îl oprește sub
(prag~$-$~2\,°C), histerezisul reducând oscilațiile rapide în jurul pragului.
Atât modul, cât și pragul sunt salvate în memoria NVS și raportate înapoi în
mesajul de senzori (\texttt{fanOn}, \texttt{fanAuto}, \texttt{fanThreshold}),
astfel încât aplicația reflectă starea reală a dispozitivului.

\subsection{Comunicație MQTT, provisioning Wi-Fi și mDNS}
\label{sec:mqtt-provisioning}

La pornire, nodul citește credențialele Wi-Fi și adresa brokerului
\gls{mqtt} din memoria NVS (\emph{Non-Volatile Storage}) folosind
biblioteca \emph{Preferences}.
Dacă rețeaua salvată nu răspunde, dispozitivul revine automat la
hotspot-ul de configurare (\texttt{SmartHome-Setup}), de unde o comandă
\texttt{wifi\_update} — primită ulterior tot prin \gls{mqtt} — poate
instala noi credențiale și reporni nodul.

Adresa brokerului poate fi specificată ca nume mDNS (\texttt{smarthome.local});
în acest caz, funcția \texttt{resolveMqttHost()} interogează serverul mDNS
local (până la cinci tentative) și returnează IP-ul corespunzător,
eliminând necesitatea configurării unui IP static pe gateway~\cite{rfc6762mdns}.

Conexiunea \gls{mqtt} este stabilită cu autentificare (utilizator și parolă)
și cu un mesaj \emph{Last Will and Testament} reținut pe topicul
\texttt{smarthome\slb{} esp32\_node\_a\slb{}status}:
dacă brokerul detectectează o deconectare bruscă, publică automat
\texttt{\{"node\_id":\allowbreak"esp32\_node\_a",\allowbreak"status":\allowbreak"offline"\}}, informând
aplicația mobilă în timp real~\cite{oasis2019mqtt}.
La conectare cu succes, nodul se abonează la
\texttt{smarthome\slb{} esp32\_node\_a\slb{}commands} și publică imediat statusul
\texttt{online}.
Publicarea datelor de senzori are loc la fiecare 5\,s pe topicul
\texttt{smarthome\slb{} esp32\_node\_a\slb{}sensors}.

% ──────────────────────────────────────────────────────────────────
\section{Firmware-ul ESP32-CAM (nod exterior)}
\label{sec:firmware-cam}

Nodul exterior este realizat pe un modul AI-Thinker ESP32-CAM echipat cu
senzorul de imagine OV2640~\cite{espressif2023esp32cam}.
Singurul senzor de mișcare din sistem — un \gls{pir} HC-SR501 — este
conectat la GPIO13~\cite{pir2019hcsr501}.
La detecția mișcării, un \gls{led} exterior conectat la GPIO12 se aprinde
timp de 30\,s (gestionat prin variabila \texttt{lastMotion} și temporizare
cu \texttt{millis()}).

\subsection{Serverul HTTP și stream-ul MJPEG}
\label{sec:mjpeg}

Camera este inițializată prin \texttt{esp\_camera\_init()}, cu rezoluție
SVGA (800$\times$600) dacă este prezentă memoria PSRAM, sau VGA în lipsa
acesteia.
Firmware-ul pornește un server HTTP pe portul 80 cu două rute:
\begin{itemize}
  \item \texttt{GET /capture} — returnează un snapshot JPEG al cadrului
        curent (\texttt{Content-Type: image/jpeg});
  \item \texttt{GET /stream} — transmite un flux \emph{Motion JPEG}
        (\texttt{multipart/x-mixed-replace}), cu o pauză de 50\,ms între
        cadre ($\approx$15–20 fps), lăsând procesorului timp și pentru
        bucla \gls{mqtt}.
\end{itemize}

Camera devine accesibilă în rețeaua locală ca
\texttt{http://smarthome-cam.local} prin înregistrarea unui serviciu mDNS
(\texttt{MDNS.begin("smarthome-cam")}), similar nodului interior.
Tot periodic (la 5\,s), nodul publică un mesaj JSON pe topicul
\texttt{smarthome/esp32\_cam\_node/sensors} conținând câmpul \texttt{motion}
(0/1), iar conexiunea \gls{mqtt} folosește același mecanism
\emph{Last Will} ca nodul interior.

% ──────────────────────────────────────────────────────────────────
\section{Gateway Raspberry Pi}
\label{sec:gateway}

Gateway-ul rulează pe un Raspberry Pi 3B+ și găzduiește toate serviciile
de infrastructură ale sistemului.

\textbf{Mosquitto} este brokerul \gls{mqtt} (portul 1883), configurat cu
autentificare prin fișier de parole.
Acesta primește publicările nodurilor ESP32, le distribuie abonaților
(\gls{api}-ului NestJS și motorului ML) și gestionează mesajele
\emph{retained} pentru statusul nodurilor.

\textbf{PostgreSQL + TimescaleDB} stochează toate citirile de senzori în
hypertabela \texttt{sensor\_readings}, optimizată pentru interogări
time-series prin indexuri pe coloana \texttt{time}~\cite{timescaledb2017,postgresql2024}.
Celelalte tabele ale schemei (\texttt{alerts}, \texttt{nodes},
\texttt{automation\_rules}, \texttt{users}, \texttt{ml\_models}) sunt
tabele PostgreSQL standard.

Fiecare componentă software este înregistrată ca serviciu \textbf{systemd},
asigurând pornirea automată la boot și repornirea în caz de eroare:
\texttt{smarthome-api.service} (NestJS),
\texttt{smarthome-ml.service} (antrenare zilnică la 03:00) și
\texttt{mosquitto.service}.

% ──────────────────────────────────────────────────────────────────
\section{API NestJS}
\label{sec:api}

\gls{api}-ul backend este implementat cu framework-ul NestJS (Node.js,
TypeScript) și rulează pe portul 3000 al gateway-ului~\cite{nestjs2024}.
TypeORM asigură maparea obiect-relațională față de PostgreSQL.

\subsection{Structura modulelor}
\label{sec:module-nestjs}

\texttt{AppModule} importă nouă module funcționale:
\texttt{AuthModule}, \texttt{UsersModule}, \texttt{SensorsModule},
\texttt{AlertsModule}, \texttt{RulesModule}, \texttt{CommandsModule},
\texttt{MqttModule}, \texttt{SimulatorModule} și \texttt{SetupModule}.
Comunicarea inter-modul se realizează prin \texttt{EventEmitter2}
(biblioteca \emph{@nestjs/event-emitter}), evitând dependențe circulare.

\subsection{Subscriber MQTT}
\label{sec:subscriber-mqtt}

\texttt{MqttService} implementează \texttt{OnModuleInit} și se conectează
la Mosquitto la pornire (dezactivat automat dacă \texttt{SIMULATE=true}).
Se abonează la trei familii de topicuri cu wildcard:
\texttt{smarthome/+/sensors} (date senzori), \texttt{smarthome/alerts}
(alerte hardware) și \texttt{smarthome/+/status} (prezența nodurilor).
\listingref{lst:mqtt-handler} prezintă logica de procesare a mesajelor.

\begin{lstlisting}[language=TypeScript,
  caption={Handler de mesaje MQTT în \texttt{MqttService}
           (\texttt{mqtt.service.ts}) — ramificarea pe tip de topic.},
  label=lst:mqtt-handler]
private async handleMessage(topic: string, payload: string) {
  const data = JSON.parse(payload);

  if (topic.endsWith('/sensors')) {
    await this.sensorsService.saveReading({
      nodeId:       data.nodeId ?? data.node_id,
      temperature:  data.temperature,
      humidity:     data.humidity,
      temperature2: data.temperature2 ?? null,
      humidity2:    data.humidity2    ?? null,
      gasLevel:     data.gasLevel,
      gasAlert:     data.gasAlert ?? false,
      light1:       data.light1   ?? null,
      light2:       data.light2   ?? null,
      lightLux:     data.lightLux ?? data.light1,
      motion:       data.motion   ?? false,
    });
  } else if (topic === 'smarthome/alerts') {
    await this.alertsService.create({
      nodeId:    data.node_id,
      alertType: data.alert_type,   // ex. "GAS_DETECTED"
      severity:  data.alert_type === 'GAS_DETECTED' ? 'critical' : 'medium',
      location:  data.location,
      details:   data.details,
    });
  } else if (topic.endsWith('/status')) {
    await this.sensorsService.updateNodeStatus(data);
  }
}
\end{lstlisting}

Publicarea comenzilor spre noduri se face prin evenimentul
\texttt{command.sent} emis de \texttt{CommandsService} și capturat de
\texttt{MqttService}, care îl publică pe topicul
\texttt{smarthome/\{nodeId\}/commands}.

\subsection{WebSocket (Socket.IO)}
\label{sec:websocket}

\texttt{SensorsGateway} declară un namespace Socket.IO \texttt{/live} și
reacționează la trei evenimente interne:
\texttt{sensor.reading} → emite \texttt{sensor\_update} (datele noi de senzori),
\texttt{alert.new} → emite \texttt{alert},
\texttt{node.status} → emite \texttt{node\_status}.
La conectarea unui client, tokenul \gls{jwt} este verificat opțional;
dacă tokenul e prezent dar invalid, clientul este deconectat imediat.

\subsection{Endpointuri REST și autentificare JWT}
\label{sec:rest}

Toate rutele protejate sunt acoperite de garda \texttt{JwtAuthGuard}.
Tabelul următor rezumă principalele endpointuri~\cite{jones2015jwt}:

\begin{table}[H]
  \centering
  \caption{Principalele endpointuri \gls{rest} ale \gls{api}-ului NestJS.}
  \label{tab:rest-endpoints}
  \begin{tabular}{@{} l l p{5.5cm} @{}}
    \toprule
    \textbf{Metodă} & \textbf{Cale} & \textbf{Descriere} \\
    \midrule
    POST & \texttt{/api/auth/login}             & Autentificare; returnează token JWT \\
    GET  & \texttt{/api/sensors/latest}         & Ultima citire per nod \\
    GET  & \texttt{/api/sensors/nodes}          & Lista nodurilor cu status online \\
    GET  & \texttt{/api/sensors/history}        & Istoric cu filtre \texttt{nodeId}, \texttt{from}, \texttt{to}, \texttt{limit} \\
    GET  & \texttt{/api/alerts}                 & Toate alertele \\
    POST & \texttt{/api/alerts/\{id\}/acknowledge} & Marchează alertă ca recunoscută \\
    POST & \texttt{/api/commands/\{nodeId\}}    & Trimite comandă actuator \\
    GET  & \texttt{/api/rules}                  & Reguli de automatizare \\
    GET  & \texttt{/api/setup/status}           & Status provisionare Wi-Fi \\
    \bottomrule
  \end{tabular}
\end{table}

Autentificarea utilizează JWT cu semnătură HS256.
Serviciul \texttt{AuthService} verifică parola prin \texttt{bcrypt.compare()}
față de hash-ul stocat în baza de date~\cite{provos2003bcrypt},
iar \texttt{JwtService.sign()} emite tokenul cu payload
\texttt{\{sub, username\}}.

\subsection{Modulul simulator}
\label{sec:simulator}

Pentru dezvoltare fără hardware fizic, variabila de mediu
\texttt{SIMULATE=true} activează \texttt{SimulatorService}.
Acesta generează la fiecare 5\,s date realiste pentru nodul interior
(\texttt{esp32\_node\_a}): temperaturi de bază de 21\,°C în living și
19\,°C în dormitor, umiditate, gaz și mișcare aleatorie, inclusiv
câmpurile \texttt{temperature2}, \texttt{humidity2}, \texttt{light1},
\texttt{light2}, asigurând paritatea cu schema reală a payload-ului.
Când \texttt{SIMULATE=false}, \texttt{MqttService} preia controlul,
iar \texttt{SimulatorService} nu pornește timer-ul.

% ──────────────────────────────────────────────────────────────────
\section{Motorul de detecție a anomaliilor cu Isolation Forest}
\label{sec:motor-ml}

Detecția anomaliilor comportamentale constituie contribuția originală a
acestei lucrări.
Motorul \gls{ml} rulează în întregime local pe Raspberry Pi, fără
dependență de cloud, în acord cu principiul \emph{local-first}
descris în \chapref{ch:metodologie}~\cite{liu2008isolation,pedregosa2011scikit}.

\subsection{Ingineria trăsăturilor (\texttt{features.py})}
\label{sec:features}

Funcția \texttt{build\_features(df)} transformă un DataFrame Pandas cu
citiri brute în matricea de trăsături necesară modelului~\cite{mckinney2010pandas}.
Cele 12 trăsături, în ordinea fixă din \texttt{FEATURE\_NAMES}, sunt
descrise în \tabref{tab:features-cod}.

\begin{table}[H]
  \centering
  \caption{Cele 12 trăsături utilizate de motorul Isolation Forest.}
  \label{tab:features-cod}
  \begin{tabular}{@{} r l p{6.5cm} @{}}
    \toprule
    \textbf{Idx} & \textbf{Trăsătură} & \textbf{Sursă / transformare} \\
    \midrule
    0  & \texttt{temperature}      & DHT11 \#1, °C \\
    1  & \texttt{humidity}         & DHT11 \#1, \%RH \\
    2  & \texttt{gas\_normalized}  & MQ-2 ADC $/ 4095$ \\
    3  & \texttt{lux\_normalized}  & $\log_{1p}(\text{light1}) / 7{,}0$ (LDR \#1, scara ADC) \\
    4  & \texttt{motion}           & \gls{pir} 0/1 (de pe ESP32-CAM, GPIO13) \\
    5  & \texttt{hour\_of\_day}    & Ora din zi, 0–23 \\
    6  & \texttt{day\_of\_week}    & 0 (luni) – 6 (duminică) \\
    7  & \texttt{is\_weekend}      & 1 dacă zi = sâmbătă sau duminică \\
    8  & \texttt{is\_night}        & 1 dacă oră $\geq 23$ sau oră $< 6$ \\
    9  & \texttt{temp\_delta\_1h}  & Diferența de temperatură față de 2 citiri anterioare \\
    10 & \texttt{motion\_count\_10}& Suma evenimentelor PIR în fereastra de 10 citiri \\
    11 & \texttt{temp\_avg}        & Media DHT11 \#1 și \#2 (fallback = DHT11 \#1) \\
    \bottomrule
  \end{tabular}
\end{table}

Trăsăturile temporale (ora, ziua, weekend, noapte) permit modelului să
distingă comportamente normale specifice perioadei (ex. mișcare
nocturnă sau temperaturi ridicate în weekend) de adevăratele anomalii.
Normalizarea \texttt{StandardScaler} (medie zero, deviație standard unu)
este aplicată în \texttt{train.py} și reutilizată în \texttt{infer.py}.

\listingref{lst:features} prezintă calculul trăsăturilor derivate.

\begin{lstlisting}[language=Python,
  caption={Calcul trăsături derivate în \texttt{features.py} —
           normalizare gaz, lumină și media DHT11.},
  label=lst:features]
df['gas_normalized'] = df['gas_level'].fillna(0) / 4095.0

# Luminozitate: LDR ADC 0-1023; log1p/7 → interval ~[0, 1]
if 'light1' in df.columns:
    raw_light = df['light1'].fillna(
        df.get('light_lux', pd.Series(0, index=df.index)).fillna(0))
else:
    raw_light = df['light_lux'].fillna(0)
df['lux_normalized'] = np.log1p(raw_light) / 7.0

# Media DHT11 #1 si #2 (daca #2 lipseste, fallback la #1)
if 'temperature2' in df.columns:
    temp2 = df['temperature2'].fillna(df['temperature'])
else:
    temp2 = df['temperature']
df['temp_avg'] = (df['temperature'] + temp2) / 2.0
\end{lstlisting}

\subsection{Antrenarea modelului (\texttt{train.py})}
\label{sec:train}

Scriptul \texttt{train.py} extrage din PostgreSQL ultimele 30 de zile de
citiri fără alertă de gaz (citirile de gaz anormal ar perturba distribuția
normală), construiește matricea de trăsături, aplică
\texttt{StandardScaler} și antrenează un \texttt{IsolationForest} cu
parametrii: \texttt{contamination=0.05} (5\% din date asumate anomale),
\texttt{n\_estimators=100}, \texttt{random\_state=42}~\cite{liu2008isolation}.
Este impus un minim de 500 de rânduri; sub acest prag, scriptul se
oprește cu mesaj de informare, fără a suprascrie modelul existent.

Modelul antrenat și scalerul sunt persistate cu \texttt{joblib}
(\texttt{isolation\_forest.pkl}, \texttt{scaler.pkl}), iar metadatele
versiunii (număr de rânduri, contaminare, lista de trăsături) sunt
înregistrate în tabela \texttt{ml\_models} din PostgreSQL.
Reantrenarea este declanșată zilnic la 03:00 de serviciul systemd
\texttt{smarthome-ml.service} prin \texttt{scheduler.py}.

\subsection{Inferența în timp real (\texttt{infer.py})}
\label{sec:infer}

Clasa \texttt{AnomalyDetector} expune metoda \texttt{score(reading: dict)}
care primește o citire nouă (dicționar Python), construiește un DataFrame
de un singur rând, aplică \texttt{build\_features()} și
\texttt{scaler.transform()}, apoi calculează
\texttt{model.\allowbreak score\_samples(\allowbreak X\_scaled)[0]}.
Un scor sub pragul \texttt{ANOMALY\_THRESHOLD = -0.1} declanșează
crearea unei alerte de tip \texttt{ML\_ANOMALY} în baza de date,
urmată de notificarea aplicației mobile prin WebSocket~\cite{chandola2009anomaly}.

\textbf{Strategia cold-start}: în primele zile de funcționare, când
baza de date nu are suficiente date pentru antrenare ($< 500$ înregistrări),
motorul revine la reguli cu prag hardcodat (gaz $> 300$,
temperatură $> 35$\,°C sau $< 5$\,°C, mișcare în intervalul 23:00–06:00),
asigurând un nivel minim de protecție până când modelul statistic poate
fi inițializat.
Funcția \texttt{is\_ready()} permite codului apelant să distingă cele
două regimuri.

% ──────────────────────────────────────────────────────────────────
\section{Fluxuri principale — diagrame de secvență}
\label{sec:secvente}

Această secțiune descrie cele trei fluxuri principale ale sistemului
prin diagrame de secvență UML.

\subsection{Fluxul de monitorizare}
\label{sec:secv-monitor}

\figref{fig:secv-monitor} ilustrează ciclul de colectare și afișare a
datelor de senzori.
La fiecare 5\,s, ESP32 citește toți senzorii și publică un mesaj JSON pe
topicul \texttt{smarthome\slb{} esp32\_node\_a\slb{}sensors}.
\texttt{MqttService} (NestJS) recepționează mesajul, apelează
\texttt{SensorsService.saveReading()} care persistă datele în
TimescaleDB, emit evenimentul intern \texttt{sensor.reading}, iar
\texttt{SensorsGateway} transmite \texttt{sensor\_update} prin
Socket.IO tuturor clienților conectați la namespace-ul \texttt{/live}.
Aplicația mobilă actualizează dashboard-ul în timp real prin
hook-ul React Query care ascultă acest eveniment.

\begin{figure}[H]
  \centering
  \includegraphics[width=\singlefigure]{\pics/fig-secventa-monitorizare}
  \caption{Diagrama de secvență pentru fluxul de monitorizare:
           ESP32 → Mosquitto → NestJS → TimescaleDB → App.}
  \label{fig:secv-monitor}
\end{figure}

\subsection{Fluxul de comandă a actuatorilor}
\label{sec:secv-comanda}

\figref{fig:secv-comanda} descrie traseul unei comenzi de la utilizator
la actuatorul fizic.
Utilizatorul apasă butonul „LED 1 ON" în ecranul Control al aplicației,
care trimite \texttt{POST /api/commands/esp32\_node\_a} cu corpul
\texttt{\{"action":"led\_on","led":0\}}.
\texttt{CommandsController} (NestJS) procesează cererea, emite evenimentul
\texttt{command.sent}, iar \texttt{MqttService} publică comanda pe topicul
\texttt{smarthome\slb{} esp32\_node\_a\slb{}commands} (QoS 1).
ESP32 recepționează mesajul în callback-ul \texttt{onMessage()},
apelează \texttt{setLed(0, true)} și aplică nivel HIGH pe GPIO25.
Întreg fluxul se desfășoară în rețeaua locală, fără ieșire pe internet.

\begin{figure}[H]
  \centering
  \includegraphics[width=\singlefigure]{\pics/fig-secventa-comanda}
  \caption{Diagrama de secvență pentru fluxul de comandă a actuatorilor:
           App → NestJS → Mosquitto → ESP32 → actuator fizic.}
  \label{fig:secv-comanda}
\end{figure}

\subsection{Fluxul de detecție a anomaliilor}
\label{sec:secv-alerta}

\figref{fig:secv-alerta} prezintă traseul complet al unei anomalii
detectate de motorul \gls{if}.
O citire nouă sosită pe topicul \texttt{smarthome/+/sensors} este salvată
în baza de date, după care \texttt{SensorsService} apelează modulul ML
(\texttt{ml.service.ts}) care invocă sincron \texttt{AnomalyDetector.score()}.
Dacă scorul returnat este sub pragul $-0{,}1$, \texttt{AlertsService}
creează o înregistrare de tip \texttt{ML\_ANOMALY} în tabela
\texttt{alerts}, emit evenimentul \texttt{alert.new}, iar
\texttt{SensorsGateway} îl transmite aplicației prin Socket.IO.
Utilizatorul primește o notificare vizuală pe ecranul Alerte.

\begin{figure}[H]
  \centering
  \includegraphics[width=\singlefigure]{\pics/fig-secventa-alerta}
  \caption{Diagrama de secvență pentru detecția anomaliei comportamentale
           cu \gls{if}: NestJS apelează \texttt{infer.py} și generează
           alertă \texttt{ML\_ANOMALY}.}
  \label{fig:secv-alerta}
\end{figure}

% ──────────────────────────────────────────────────────────────────
\section{Aplicația mobilă}
\label{sec:mobile}

Interfața utilizatorului este implementată ca aplicație React Native 0.81
compilată cu Expo SDK 54, targetând iOS și Android~\cite{reactnative2024,expo2024}.
Starea globală este gestionată cu \textbf{Zustand}, apelurile \gls{rest}
cu \textbf{React Query}, iar datele în timp real prin \textbf{Socket.IO}
conectat la namespace-ul \texttt{/live} al \gls{api}-ului.
Autentificarea utilizează token \gls{jwt} stocat local.

\subsection{Discovery-ul gateway-ului}
\label{sec:discovery}

La prima pornire, aplicația încearcă să contacteze gateway-ul în
ordinea: variabila de mediu configurată la build-time,
\texttt{http://smarthome.local:3000} (mDNS) sau adresa fixă a
hotspot-ului de configurare (\texttt{192.168.4.1}).
Această strategie elimină nevoia configurării manuale a IP-ului în
scenariile tipice de rețea locală, susținând obiectivul de funcționare
„plug and play". Implementarea, prezentată în \listingref{lst:discovery},
sondează \emph{în paralel} toate adresele candidate și o reține pe prima
care răspunde la endpointul public \texttt{/api/setup/status}; la o eroare
de rețea adresa este invalidată, astfel încât aplicația redescoperă
gateway-ul atunci când telefonul trece, de exemplu, de pe hotspot pe
rețeaua Wi-Fi a casei.

\begin{lstlisting}[language=TypeScript,
  caption={Descoperirea automată a gateway-ului în rețea (\texttt{services/discovery.ts}).},
  label=lst:discovery]
const CANDIDATES = [
  API_HOST,                       // variabila de mediu / fallback
  'http://smarthome.local:3000',  // mDNS - orice retea
  HOTSPOT_HOST,                   // hotspot RPi - 192.168.4.1
];

// sondeaza in paralel; prima adresa care raspunde "castiga"
export async function resolveHost(force = false): Promise<string> {
  if (resolvedHost && !force) return resolvedHost;
  return Promise.any(CANDIDATES.map(probe))
    .then((host) => { resolvedHost = host; return host; })
    .catch(() => API_HOST);
}

async function probe(host: string): Promise<string> {
  // /api/setup/status e public si raspunde doar de pe RPi
  await axios.get(`${host}/api/setup/status`, { timeout: 4000 });
  return host;
}
\end{lstlisting}

\subsection{Recepția datelor în timp real (Socket.IO)}
\label{sec:mobile-ws}

Actualizarea live a interfeței nu se bazează pe interogare repetată
(\emph{polling}), ci pe o conexiune \textbf{Socket.IO} către namespace-ul
\texttt{/live} al \gls{api}-ului, autentificată cu același token \gls{jwt}.
La fiecare eveniment primit de la server, handler-ul actualizează direct
store-ul Zustand, iar componentele abonate se redesenează automat
(\listingref{lst:ws-client}). Reconectarea este automată, cu un număr
nelimitat de încercări — esențial pentru un dispozitiv mobil care își
schimbă frecvent starea rețelei.

\begin{lstlisting}[language=TypeScript,
  caption={Clientul Socket.IO: evenimente live $\rightarrow$ store global (\texttt{services/socket.ts}).},
  label=lst:ws-client]
socket = io(`${host}/live`, {
  auth: { token },
  transports: ['websocket'],
  reconnectionDelay: 3000,
  reconnectionAttempts: Infinity,
});

// citire noua de senzor -> actualizeaza store-ul -> UI live
socket.on('sensor_update', (reading: SensorReading) => {
  useAppStore.getState().setLatestReading(reading);
});
socket.on('alert', (alert: AlertItem) => {
  useAppStore.getState().addAlert(alert);          // badge + ecran Alerte
});
socket.on('node_status', (s: { nodeId: string; online: boolean }) => {
  useAppStore.getState().setNodeStatus(s.nodeId, s.online);
});
\end{lstlisting}

\subsection{Trimiterea comenzilor către actuatori}
\label{sec:mobile-cmd}

Comenzile către actuatori sunt expuse printr-o singură funcție generică
\texttt{api.\allowbreak commands.\allowbreak send(nodeId,\allowbreak action,\allowbreak extras)}, care construiește
corpul \gls{json} al cererii \gls{rest}; \texttt{extras} permite atașarea
parametrilor specifici (de exemplu indexul \gls{led}-ului sau unghiul
servomotorului). În ecranul de control, fiecare acțiune este înfășurată
într-o mutație React Query care aplică o \emph{actualizare optimistă} a
stării vizuale, astfel încât comutatorul reacționează instantaneu
(\listingref{lst:mobile-cmd}).

\begin{lstlisting}[language=TypeScript,
  caption={Trimiterea unei comenzi de actuator cu actualizare optimistă (\texttt{services/api.ts} + \texttt{ControlScreen.tsx}).},
  label=lst:mobile-cmd]
// services/api.ts - o singura functie pentru toate comenzile
send: (nodeId, action, extras) =>
  client.post(`/commands/${nodeId}`, { action, ...extras })
        .then((r) => r.data),

// ControlScreen.tsx - buton LED cu actualizare optimista
const mutation = useMutation({
  mutationFn: ({ action, led }) => api.commands.send(nodeId, action, { led }),
  onSuccess: (_, vars) =>
    setLeds((prev) => ({ ...prev, [vars.led]: vars.action === 'led_on' })),
});
// la apasarea butonului:
mutation.mutate({ action: on ? 'led_off' : 'led_on', led: led.index });
\end{lstlisting}

\subsection{Ecranele aplicației}
\label{sec:ecrane}

Aplicația este organizată în patru tab-uri principale accesibile după
autentificare, plus un ecran de setup:

\textbf{Dashboard} afișează un „Hero card" cu temperatura medie,
umiditatea, iluminarea și nivelul de gaz agregate de pe toate nodurile,
urmat de tile-uri individuale pentru fiecare nod.
Ecranul prezintă distinct datele celor doi senzori \gls{dht}11
(\texttt{temperature}/\texttt{humidity} pentru living și
\texttt{temperature2}/\texttt{humidity2} pentru dormitor)
și valorile brute ADC ale celor două \gls{ldr}-uri (\texttt{light1},
\texttt{light2}).
Sparkline-uri animate vizualizează tendința recentă a fiecărei valori,
iar un indicator „scor de siguranță" sintetizează starea generală.

\textbf{Control} permite acționarea tuturor actuatorilor:
\begin{itemize}
  \item Trei butoane individuale pentru \gls{led}-urile din living
        (GPIO25), dormitor (GPIO26) și hol (GPIO27);
  \item Comutator pentru ventilator (GPIO33 prin tranzistor);
  \item Trei preset-uri pentru servo (draperii): Deschis (0°), 50\%
        (90°), Închis (180°);
  \item Buton de test buzzer și buton „Oprire totală" (trimite comanda
        \texttt{all\_off});
  \item Scene predefinite: Acasă, Plec, Noapte, Cinema;
  \item Vizualizare cameră: snapshot (polling la 5\,s via
        \texttt{GET /capture}) și stream live MJPEG redat în
        \texttt{WebView} via \texttt{GET /stream}.
\end{itemize}

\textbf{Istoric} prezintă grafice de tendință pentru parametrii
selectați (temperatură, umiditate, gaz, lumină) pe intervale de
timp configurabile, folosind date din \texttt{GET /api/sensors/history}.

\textbf{Alerte} listează cronologic toate alertele generate de sistem
(\texttt{GAS\_DETECTED}, \texttt{MOTION\_NIGHT}, \texttt{ALERT\_HEAT},
\texttt{ALERT\_COLD}, \texttt{ML\_ANOMALY}), cu indicatori de severitate
și posibilitatea de a marca o alertă ca recunoscută
(\texttt{POST /api/alerts/\{id\}/acknowledge}).

\figref{fig:mobile} prezintă capturi reale din aplicație, ilustrând
limbajul vizual \emph{glassmorphism} adoptat: suprafețe translucide cu
estompare (\textit{blur}), accente vii pe gradient (indigo–cyan),
tipografie Space~Grotesk pentru valori numerice și Inter pentru text,
respectând contrastul \gls{ui} conform recomandărilor de accesibilitate.

\begin{figure}[H]
  \centering
  \subfloat[Dashboard]{\includegraphics[width=0.235\textwidth]{\pics/fig-mobile-dashboard}}\hfill
  \subfloat[Alerte]{\includegraphics[width=0.235\textwidth]{\pics/fig-mobile-alerte}}\hfill
  \subfloat[Control (sus)]{\includegraphics[width=0.235\textwidth]{\pics/fig-mobile-control1}}\hfill
  \subfloat[Control (jos)]{\includegraphics[width=0.235\textwidth]{\pics/fig-mobile-control2}}
  \caption{Capturi reale din aplicația mobilă: (a)~Dashboard — scor de
           siguranță, medii ambientale și tile-uri pe camere; (b)~Alerte —
           cronologie cu severitate și stare de confirmare; (c)–(d)~Control —
           cameră live, scene, draperii (servo), \gls{led}-uri și ventilator,
           cu indicator „nod offline" când nodul nu este conectat.}
  \label{fig:mobile}
\end{figure}

% ──────────────────────────────────────────────────────────────────
\section{Sumar}
\label{sec:implementare-sumar}

Capitolul a prezentat implementarea concretă a fiecărui strat al
sistemului smart home local-first.
Firmware-ul nodului interior (ESP32 WROOM-32) gestionează cei unsprezece
pini activi, publică un payload JSON standardizat la 5\,s și asigură
reconectarea automată la broker cu mecanisme \emph{Last Will} și
provisioning \gls{wifi} prin NVS.
ESP32-CAM adaugă capacitate video (stream MJPEG, snapshot JPEG) și
detecție de mișcare prin \gls{pir} GPIO13.
Gateway-ul Raspberry Pi orchestrează Mosquitto, PostgreSQL/TimescaleDB
și serviciile systemd.
\gls{api}-ul NestJS leagă toate componentele prin module clare (MQTT,
WebSocket, \gls{rest}, auth \gls{jwt}).
Contribuția originală — motorul Isolation Forest cu 12 trăsături —
rulează local și detectează anomalii comportamentale în timp real,
cu strategie cold-start pentru fazele inițiale de colectare a datelor.
Aplicația mobilă React Native/Expo oferă vizualizare live, control
complet al actuatorilor și acces la camera exterioară, toate fără
dependență de cloud.

% cap5_rezultate.tex
\chapter{Rezultate și evaluare}
\label{ch:rezultate}

Prezentul capitol descrie metodologia de evaluare adoptată pentru prototipul de sistem smart home
local-first și prezintă rezultatele obținute în faza curentă de dezvoltare. Sistemul se află la
stadiul de prototip funcțional asamblat pe breadboard; modelul de \gls{ml} bazat pe
\gls{if} nu a fost încă antrenat pe date reale de lungă durată, iar suita de teste automate
este în curs de implementare. Prin urmare, orice cifră numerică prezentată în acest capitol este
marcată explicit ca \textbf{valoare preliminară / ilustrativă} și reflectă fie măsurători punctuale
obținute în sesiuni de testare manuală, fie estimări bazate pe specificațiile tehnice ale
componentelor utilizate. Acestea urmează să fie validate statistic pe o perioadă mai lungă de
operare.

% ──────────────────────────────────────────────────────────────────────
\section{Mediul și metodologia de testare}
\label{sec:mediu-testare}
% ──────────────────────────────────────────────────────────────────────

\subsection{Configurația hardware de testare}

Întregul prototip a fost asamblat pe breadboard, fără circuit imprimat dedicat, ceea ce introduce
variabilitate suplimentară în măsurătorile de tensiune și timpi de răspuns. Componentele principale
sunt:

\begin{itemize}
  \item un nod interior bazat pe ESP32 WROOM-32, cu senzorii \gls{dht}11 (×2), MQ-2,
        \gls{ldr} (×2), buzzer, servomotor SG90 și trei \gls{led}-uri;
  \item un nod exterior ESP32-CAM (AI-Thinker, OV2640), cu senzor \gls{pir} pe GPIO13;
  \item un gateway Raspberry Pi 3B+ rulând Mosquitto (port 1883), PostgreSQL + TimescaleDB,
        \gls{api}-ul NestJS (port 3000) și motorul \gls{ml} Python.
\end{itemize}

Rețeaua locală de testare este un router Wi-Fi de uz casnic; descoperirea serviciilor se realizează
prin \gls{mdns} (\texttt{smarthome.local}, \texttt{smarthome-cam.local}).

\subsection{Modul simulator}

Deoarece senzorii fizici produc date cu rată de 5 secunde și necesită prezența fizică pentru
stimularea unor scenarii (de exemplu gaz detectat), \gls{api}-ul NestJS include un modul
\texttt{simulator/} activat prin variabila de mediu \texttt{SIMULATE=true}. Acesta generează
citiri sintetice cu variații realiste ale tuturor câmpurilor (inclusiv \texttt{temperature2},
\texttt{humidity2}, \texttt{light1}, \texttt{light2}), permițând testarea fluxurilor de date,
a persistenței în baza de date și a comunicării \gls{mqtt} fără hardware conectat.

\subsection{Piramida de testare propusă}

Metodologia de testare urmează structura clasică în trei niveluri, ilustrată în
\figref{fig:piramida}.

\begin{figure}[H]
  \centering
  \includegraphics[width=\singlefigure]{\pics/fig-piramida-testare}
  \caption{Piramida de testare propusă pentru sistem: teste unitare la bază, teste de integrare la
    nivel intermediar și teste end-to-end la vârf. Lateral — testarea manuală pe hardware.}
  \label{fig:piramida}
\end{figure}

\textbf{Nivelul 1 — Teste unitare.}
Serviciile \gls{api}-ului NestJS (autentificare, publicare comenzi, procesare alertă) și modulul
\texttt{features.py} (calculul celor 12 trăsături, normalizarea cu StandardScaler) sunt vizate
pentru testare cu Jest, respectiv \texttt{pytest}. Aceste teste verifică logica internă a fiecărei
componente în izolare, cu dependențe substituibile prin mock-uri.

\textbf{Nivelul 2 — Teste de integrare.}
Scenariile de integrare vizează lanțul complet: \gls{api} NestJS $\leftrightarrow$ broker
Mosquitto $\leftrightarrow$ baza de date PostgreSQL/TimescaleDB. Se verifică, de exemplu,
că o citire publicată pe topicul \texttt{smarthome/esp32\_node\_a/sensors} este persistată corect
în hypertable-ul \texttt{sensor\_readings} și propagată prin WebSocket pe canalul
\texttt{/live}.

\textbf{Nivelul 3 — Teste end-to-end.}
Scenariile end-to-end acoperă lanțul integral firmware $\leftrightarrow$ broker $\leftrightarrow$
\gls{api} $\leftrightarrow$ aplicație mobilă. Se verifică că o comandă emisă din interfața mobilă
ajunge la actuatorul fizic în sub-secundă și că o alertă de gaz generată de firmware apare pe
ecranul utilizatorului.

\textbf{Testare manuală pe hardware.}
Paralel cu testele automate (în curs de implementare), fiecare funcționalitate a fost verificată
manual: conectarea nodului la rețea, publicarea datelor \gls{mqtt}, executarea comenzilor de
actuator, generarea alertei de gaz și vizualizarea fluxului video de la camera ESP32-CAM.

% ──────────────────────────────────────────────────────────────────────
\section{Verificarea funcțională a sistemului}
\label{sec:verificare-functionala}
% ──────────────────────────────────────────────────────────────────────

Tabelul~\tabref{tab:verif-functionala} centralizează funcționalitățile principale ale sistemului,
metoda de verificare aplicată și statusul curent. Verificarea „manuală confirmat" înseamnă că
funcționalitatea a fost testată în sesiuni de lucru pe prototipul fizic sau cu simulatorul activat.

\begin{table}[H]
  \centering
  \caption{Verificarea funcțională a componentelor sistemului.}
  \label{tab:verif-functionala}
  \small
  \begin{tabular}{p{4.8cm} p{5.4cm} p{2.6cm}}
    \toprule
    \textbf{Funcționalitate} & \textbf{Metodă de verificare} & \textbf{Status} \\
    \midrule
    Conectare ESP32 la broker \gls{mqtt} (mDNS) &
      Monitorizare log Mosquitto; LED de stare GPIO25 &
      Confirmat manual \\
    \addlinespace
    Publicare citiri senzori la 5\,s pe topicul \texttt{sensors} &
      Subscriber MQTT extern (\texttt{mosquitto\_sub}); vizualizare dashboard &
      Confirmat manual \\
    \addlinespace
    Persistare în TimescaleDB (hypertable) &
      Interogare SQL pe \texttt{sensor\_readings}; endpointul \texttt{/api/sensors/history} &
      Confirmat manual \\
    \addlinespace
    Comanda \texttt{led\_on/off/toggle} (LED 0, 1, 2) &
      Observare vizuală; log firmware serial &
      Confirmat manual \\
    \addlinespace
    Comanda \texttt{fan\_on/off} (ventilator prin tranzistor) &
      Observare fizică + log serial &
      Confirmat manual \\
    \addlinespace
    Comanda \texttt{servo\_move} (0°–180°, draperii) &
      Observare mișcare servomotor SG90 &
      Confirmat manual \\
    \addlinespace
    Comanda \texttt{buzzer\_beep} / \texttt{all\_off} &
      Observare sonoră + log serial &
      Confirmat manual \\
    \addlinespace
    Alertă gaz (\texttt{gasLevel > 350}) → \texttt{smarthome/alerts} + beep &
      Injectare nivel gaz simulat; observare alertă în aplicație &
      Confirmat manual \\
    \addlinespace
    Propagare alertă prin WebSocket → aplicație mobilă &
      Vizualizare tab Alerte în timp real &
      Confirmat manual \\
    \addlinespace
    Stare nod (Last Will \gls{mqtt}, \texttt{status} retained) &
      Oprire bruscă ESP32; verificare flag \texttt{online=false} &
      Confirmat manual \\
    \addlinespace
    Stream video ESP32-CAM (\gls{mjpeg}) + snapshot &
      Acces \url{http://smarthome-cam.local/stream} din WebView &
      Confirmat manual \\
    \addlinespace
    Autentificare \gls{jwt} + bcrypt &
      Testare endpoint \texttt{/api/auth/login}; token valid/expirat &
      Confirmat manual \\
    \addlinespace
    Provisioning Wi-Fi (hotspot \texttt{SmartHome-Setup}) &
      Conectare la AP fallback; trimitere comandă \texttt{wifi\_update} &
      Confirmat manual \\
    \addlinespace
    Inferență \gls{ml} (cold-start, reguli prag) &
      Injectare valori extreme; verificare alertă \texttt{ML\_ANOMALY} &
      Confirmat manual \\
    \addlinespace
    Simulator \texttt{SIMULATE=true} — date mock &
      Pornire \gls{api} în mod simulator; verificare dashboard &
      Confirmat manual \\
    \addlinespace
    Teste automate (Jest / pytest) &
      Suită în curs de implementare &
      \textit{În lucru} \\
    \bottomrule
  \end{tabular}
\end{table}

\subsection{Fluxul de comandă}

Fluxul de comandă pornește de la interacțiunea utilizatorului în aplicația mobilă (de exemplu,
apăsarea butonului „LED 1 ON") și traversează: \gls{api} NestJS → publicare \gls{mqtt} pe topicul
\texttt{smarthome/esp32\_node\_a/commands} → recepție de către firmware → execuție \texttt{setLed(0,true)}.
Confirmarea nu este transmisă explicit înapoi (firmware-ul nu publică un ACK), dar starea
actuală se reflectă în citirea imediat următoare dacă un senzor asociat este prezent (de exemplu,
schimbarea iluminării detectată de \gls{ldr}).

\subsection{Fluxul de alertă de gaz}

Detecția gazului funcționează pe front crescător: firmware-ul evaluează la fiecare citire dacă
\texttt{gasLevel} (valoare \gls{adc} pe 10 biți, 0–1023) depășește pragul de 350. La depășire,
dacă alerta nu era deja activă (\texttt{gasAlert == false}), firmware-ul publică imediat pe
\texttt{smarthome/alerts}, activează buzzer-ul local (beep autonom, independent de rețea) și
setează \texttt{gasAlert = true}. Alerta dispare (front descrescător) când nivelul scade sub 350.
\gls{api}-ul persistă evenimentul ca \texttt{GAS\_DETECTED} și îl propagă prin WebSocket.

% ──────────────────────────────────────────────────────────────────────
\section{Performanță operațională}
\label{sec:performanta}
% ──────────────────────────────────────────────────────────────────────

\textit{Notă: valorile din această secțiune sunt \textbf{preliminare / ilustrative}, obținute
din sesiuni de testare manuală punctuală pe prototipul fizic sau estimate pe baza specificațiilor
tehnice. Ele nu au caracter statistic și urmează să fie validate prin măsurători sistematice pe
o perioadă extinsă.}

\subsection{Latența comandă–actuator}

Latența este definită ca intervalul dintre momentul în care aplicația mobilă emite cererea
\texttt{POST /api/commands/\{nodeId\}} și momentul în care firmware-ul execută acțiunea (observabil
prin \texttt{Serial.print} pe UART). Factorii care o influențează includ latența \gls{wifi} în
\gls{lan}, timpul de procesare al broker-ului Mosquitto și bucla principală a firmware-ului
Arduino.

\subsection{Frecvența de publicare senzori}

Intervalul de publicare este configurat în firmware la 5 secunde (\texttt{PUBLISH\_INTERVAL = 5000\,ms}).
Citirea senzorilor \gls{dht}11 durează tipic 250\,ms per senzor conform specificațiilor
producătorului~\cite{aosong2019dht11}, deci ciclul complet (2 × DHT11 + MQ-2 + 2 × \gls{ldr})
se încadrează confortabil sub pragul de 5 secunde.

\begin{table}[H]
  \centering
  \caption{Indicatori de performanță operațională — valori \textbf{preliminare / ilustrative}.}
  \label{tab:performanta}
  \small
  \begin{tabular}{p{5.0cm} p{3.2cm} p{4.4cm}}
    \toprule
    \textbf{Indicator} & \textbf{Valoare ilustrativă} & \textbf{Observații} \\
    \midrule
    Latență comandă→actuator (\gls{lan}) &
      $\approx 120$–$250\,\text{ms}$ &
      Măsurată manual (timestamp log serial); variabilă în funcție de încărcare Wi-Fi \\
    \addlinespace
    Interval publicare citiri senzori &
      $5\,\text{s}$ (configurat) &
      Parametru constant în firmware; citire DHT11 $\approx 250\,\text{ms}$/senzor \\
    \addlinespace
    Timp răspuns \gls{api} (\texttt{/api/sensors/latest}) &
      $< 50\,\text{ms}$ &
      Măsurare locală cu \texttt{curl}; fără latență de rețea externă \\
    \addlinespace
    Consum \gls{cpu} \gls{rpi} (idle, toate serviciile active) &
      $\approx 15$–$25\%$ &
      Estimare \texttt{htop}; vârfuri la reantrenare ML (03:00) \\
    \addlinespace
    Consum \gls{ram} \gls{rpi} (toate serviciile active) &
      $\approx 450$–$550\,\text{MB}$ &
      PostgreSQL + NestJS + Mosquitto + Python; din totalul de 1\,GB \\
    \addlinespace
    Timp inferență \gls{ml} per citire (\texttt{infer.py}) &
      $< 5\,\text{ms}$ &
      Estimare pe baza caracteristicilor \gls{if}; neglijabil față de intervalul de 5\,s \\
    \addlinespace
    Latență stream video (\gls{mjpeg}) în \gls{lan} &
      $\approx 300$–$800\,\text{ms}$ &
      Observare subiectivă în WebView; dependentă de rezoluție OV2640 \\
    \bottomrule
  \end{tabular}
\end{table}

Rezultatele preliminare indică faptul că sistemul este operabil pe hardware de uz general (\gls{rpi}
3B+) fără a atinge limitele de resurse în condiții normale de funcționare. Vârfurile de \gls{cpu}
la reantrenarea zilnică a modelului \gls{ml} (programată la ora 03:00 prin \texttt{scheduler.py})
nu afectează funcționarea în timp real datorită programării în afara intervalului de utilizare.

% ──────────────────────────────────────────────────────────────────────
\section{Evaluarea detecției de anomalii}
\label{sec:eval-anomalii}
% ──────────────────────────────────────────────────────────────────────

Detecția de anomalii comportamentale cu \gls{if} reprezintă contribuția originală a acestei
lucrări. Aceasta este prezentată detaliat din perspectivă algoritmică în \chapref{ch:metodologie}.
Prezenta secțiune descrie \textbf{protocolul de evaluare propus} și discută parametrii cheie ai
modelului, cu mențiunea că evaluarea cantitativă completă urmează a fi realizată după acumularea
unui set de date real de cel puțin 30 de zile de operare continuă.

\subsection{Protocolul de evaluare propus}

Evaluarea modelului \gls{if} se va realiza prin \textbf{injectare de anomalii sintetice} în date
istorice reale. Procedura propusă cuprinde pașii următori:

\begin{enumerate}
  \item \textbf{Colectarea datelor de referință.} Se colectează minim 30 de zile de citiri
    normale în condiții de viață reală (\texttt{sensor\_readings}), cu o rată de 5 secunde,
    reprezentând un set de referință de cel puțin $\sim$500.000 de rânduri.

  \item \textbf{Antrenarea modelului de referință.} Se antrenează modelul cu parametrii actuali:
    \texttt{IsolationForest(contamination=0.05, n\_estimators=100)} pe datele de referință,
    cu normalizare StandardScaler pe cele 12 trăsături din \texttt{features.py}.

  \item \textbf{Generarea anomaliilor sintetice.} Se injectează perturbații controlate în setul
    de test:
    \begin{itemize}
      \item spike-uri de temperatură ($+15\,°C$ față de medie, durată 1–5 minute),
      \item nivel de gaz ridicat artificial (\texttt{gasLevel > 600}, necorelat cu temperatura),
      \item mișcare detectată în interval de noapte (23:00–06:00) în absența altor activități,
      \item combinații atipice de trăsături (de exemplu umiditate 95\% corelată cu temperatură scăzută).
    \end{itemize}

  \item \textbf{Evaluarea cu metrici standard.} Se calculează precizia ($P$), recall-ul ($R$),
    scorul $F_1$ și curba \gls{roc} cu aria \gls{auc}-ROC~\cite{fawcett2006roc,davis2006prcurve}:
    \[
      F_1 = \frac{2 \cdot P \cdot R}{P + R}, \qquad
      \text{AUC-ROC} = \int_0^1 \text{TPR}(u)\,\mathrm{d}u.
    \]
    Alegerea pragului optimal se realizează prin analiza curbei Precizie-Recall~\cite{davis2006prcurve},
    preferată față de curba ROC în seturile dezechilibrate (anomaliile sunt rare).

  \item \textbf{Sensibilitate la prag.} Se variază \texttt{ANOMALY\_THRESHOLD} în intervalul
    $[-0.3, 0.0]$ pentru a cuantifica compromisul fals-pozitive vs.\ fals-negative.
\end{enumerate}

\subsection{Parametrii actuali ai modelului}

Modelul curent, descris în \texttt{gateway/ml/train.py}, folosește:
\begin{itemize}
  \item \textbf{contamination = 0.05}: se presupune că 5\% din datele de antrenament pot fi
    anomalii; aceasta determină pragul intern al arborelui de izolare~\cite{liu2008isolation};
  \item \textbf{ANOMALY\_THRESHOLD = $-0.1$}: scorul \texttt{score\_samples} sub $-0.1$
    declanșează o alertă \texttt{ML\_ANOMALY}; valori mai negative corespund unor puncte mai
    ușor de izolat (mai anomale);
  \item \textbf{n\_estimators = 100}: 100 de arbori de izolare, suficient pentru stabilizarea
    estimărilor conform literaturii~\cite{liu2008isolation,liu2012isolationtkdd};
  \item \textbf{reantrenare zilnică la 03:00} prin \texttt{scheduler.py} / serviciu
    \texttt{systemd}, pe ultimele 30 de zile de date.
\end{itemize}

Alegerea pragului $-0.1$ este un punct de start conservator: permite detectarea anomaliilor
moderate fără a inunda utilizatorul cu false pozitive. Calibrarea finală va depinde de datele
reale și de preferința utilizatorului față de compromisul precizie–recall.

\subsection{Mecanism de cold-start}

Deoarece modelul \gls{if} necesită minim $\sim$500 rânduri de antrenament (aprox. 14 zile la
5 secunde/citire cu pierderi de conectivitate incluse), sistemul funcționează în mod
\textbf{cold-start} în primele săptămâni, folosind reguli prag hardcodate: gaz $>$ 300,
temperatură $> 35\,°C$ sau $< 5\,°C$, mișcare în intervalul 23:00–06:00. Odată generat primul
model valid, \texttt{infer.py} comută automat pe inferența bazată pe scor.

\subsection{Rezultate ilustrative}

Figura~\figref{fig:eval} prezintă un exemplu ilustrativ al comportamentului așteptat al
detectorului de anomalii pe un interval de 24 de ore, generat cu date sintetice.

\begin{figure}[H]
  \centering
  \includegraphics[width=\singlefigure]{\pics/fig-eval-anomalii}
  \caption{Exemplu ilustrativ al scorului de anomalie pe 24\,h (date sintetice / preliminare):
    serie temporală a scorului \gls{if}, prag la $-0.1$ și puncte detectate ca anomalii (roșu).
    \textbf{Notă: datele sunt ilustrative și nu provin din operare reală de durată.}}
  \label{fig:eval}
\end{figure}

Graficul ilustrează mecanismul de detectare: scorul se menține în intervalul $[-0.05, 0.05]$
pentru comportamente normale, coboară sub $-0.1$ la injectarea unor spike-uri sintetice de
temperatură și gaz, și revine la normal odată ce perturbația dispare. Pe datele sintetice folosite
în faza de dezvoltare, parametrii actuali (\texttt{contamination=0.05}, prag $-0.1$) au produs
rezultate orientative de $F_1 \approx 0.78$ — valoare \textbf{strict ilustrativă}, obținută
pe un set mic de anomalii injectate manual, fără semnificație statistică.

% ──────────────────────────────────────────────────────────────────────
\section{Discuție: avantaje confirmate și limitări actuale}
\label{sec:discutie}
% ──────────────────────────────────────────────────────────────────────

\subsection{Avantaje confirmate ale abordării}

\textbf{Funcționare local-first.}
Verificarea funcțională a confirmat că toate fluxurile principale (monitorizare, control, alertă,
inferență \gls{ml}) funcționează integral în \gls{lan}, fără dependențe de servicii externe.
Deconectarea de la internet nu afectează niciuna dintre funcționalitățile de bază, ceea ce
aliniază sistemul cu principiile local-first~\cite{kleppmann2019localfirst} și adresează
preocupările de confidențialitate documentate în literatura \gls{gdpr}~\cite{zheng2018smarthomeprivacy,gdpr2016}.

\textbf{Cost redus al platformei hardware.}
Arhitectura bazată pe ESP32 (~3–5 EUR/unitate) și \gls{rpi} 3B+ (~35 EUR) demonstrează că
un sistem smart home cu \gls{ml} integrat este realizabil la un cost total al componentelor
sub 60 EUR, față de soluțiile comerciale care implică abonamente lunare și dependențe de cloud.

\textbf{Inferență \gls{ml} la margine (edge).}
Execuția modelului \gls{if} direct pe \gls{rpi}, cu timpi de inferență sub 5\,ms per citire,
confirmă fezabilitatea abordării edge-ML pentru scenarii \gls{iot} rezidențiale~\cite{shi2016edge,premsankar2018edgeiot}.
Modelul nesupervizat nu necesită date etichetate, avantaj esențial pentru utilizatorul rezidențial
care nu poate produce etichete manuale~\cite{chandola2009anomaly,cook2020anomalyiot}.

\textbf{Modularitate și extensibilitate.}
Separarea clară între firmware (Arduino C++), gateway (NestJS + Python) și aplicație mobilă
(React Native/Expo) facilitează adăugarea de noi noduri, senzori sau reguli de automatizare fără
modificarea celorlalte straturi.

\subsection{Limitări actuale}

\textbf{Absența datelor reale de lungă durată.}
Cea mai semnificativă limitare actuală este că modelul \gls{if} nu a fost antrenat pe date
reale provenite din operarea continuă a sistemului. Performanța sa reală — în termeni de
precizie, recall și $F_1$ — nu poate fi afirmată cu certitudine și urmează să fie măsurată
după cel puțin 30–60 de zile de colectare continuă.

\textbf{Absența \gls{tls} pe \gls{mqtt}.}
Comunicarea \gls{mqtt} curentă se realizează pe portul 1883 (plaintext), cu autentificare
simplă prin nume de utilizator și parolă. Deși acest lucru este acceptabil într-o rețea locală
controlată, adăugarea \gls{tls} (portul 8883) este planificată pentru versiunile viitoare,
în conformitate cu recomandările de securitate pentru dispozitive \gls{iot}~\cite{alrawi2019sokhome}.

\textbf{Absența testelor automate.}
Suita de teste Jest (pentru NestJS) și pytest (pentru \texttt{features.py}) nu este încă
implementată. Aceasta limitează capacitatea de regresie și face verificarea funcționalităților
dependentă de testarea manuală.

\textbf{Prototip pe breadboard.}
Asamblarea pe breadboard introduce conexiuni nesigure și variabilitate în măsurătorile analogice
(\gls{adc}), în special pentru senzorul MQ-2 și \gls{ldr}-uri. Migrarea pe un circuit imprimat
dedicat (PCB) sau cel puțin pe o placă de prototipare permanentă (perfboard) ar îmbunătăți
repetabilitatea măsurătorilor.

\textbf{Scalabilitate — nod unic interior.}
Topologia curentă include un singur nod interior. Extinderea la mai multe noduri (camere
diferite) ar necesita ajustarea modelului \gls{ml} pentru a trata fiecare locație independent
sau a construi un model agregat.

\textbf{Calibrarea senzorului MQ-2.}
Senzorul MQ-2~\cite{hanwei2020mq2} necesită o perioadă de preîncălzire (24–48 ore) și o
calibrare în aer curat (\textit{R\textsubscript{0}}) pentru a produce valori absolute de
concentrație (ppm). În prototipul curent se folosesc valori ADC brute (0–1023), suficiente
pentru detecția unui prag de alertă dar necorespunzătoare pentru cuantificarea precisă a
concentrației de gaz.

% ──────────────────────────────────────────────────────────────────────
\section{Sinteză}
\label{sec:sinteza-rezultate}
% ──────────────────────────────────────────────────────────────────────

Capitolul a prezentat metodologia de evaluare adoptată pentru prototipul de sistem smart home
local-first și a raportat starea curentă a verificărilor. Toate funcționalitățile principale
— de la publicarea datelor \gls{mqtt} și persistarea în TimescaleDB, la controlul actuatoarelor
și propagarea alertelor prin WebSocket — au fost confirmate prin testare manuală. Indicatorii
de performanță operațională (latență, consum resurse pe \gls{rpi}, timp inferență \gls{ml})
au fost prezentați ca valori ilustrative/preliminare, obținute în sesiuni punctuale de testare.

Protocolul de evaluare a detecției de anomalii cu \gls{if} a fost definit complet, inclusiv
metrica de injectare de anomalii sintetice, metricile $F_1$ și AUC-ROC și analiza sensibilității
la prag. Evaluarea cantitativă definitivă va deveni posibilă după acumularea datelor de operare
reală. Limitările principale identificate — absența \gls{tls} pe \gls{mqtt}, lipsa testelor
automate și dependența de date de durată pentru model — constituie direcții concrete de
dezvoltare ulterioară, discutate în \chapref{ch:concluzii}.

% cap6_concluzii.tex
\chapter{Concluzii și dezvoltări viitoare}
\label{ch:concluzii}

%=============================================================================
\section{Concluzii generale}
%=============================================================================

Lucrarea de față a propus și a implementat un sistem de casă inteligentă
\emph{local-first}, funcțional integral în rețeaua locală, fără dependență de
servicii cloud externe~\cite{kleppmann2019localfirst}.  Obiectivul central
formulat în \chapref{ch:introducere} — de a demonstra că securitatea,
confidențialitatea și autonomia unui mediu rezidențial pot fi asigurate prin
mijloace accesibile, rulate pe hardware de consum — a fost atins prin
construirea unui sistem stratificat, de la senzorii de câmp până la interfața
mobilă a utilizatorului.

Arhitectura realizată cuprinde cinci straturi funcționale: (1)~stratul
senzori și actuatori, compus din 11 piese hardware; (2)~stratul noduri
\gls{iot} pe bază de \mbox{ESP32}; (3)~gateway-ul Raspberry Pi
(\gls{rpi}), care găzduiește broker-ul \gls{mqtt} Mosquitto, baza de date
PostgreSQL/TimescaleDB și motorul de \gls{ml}; (4)~\gls{api}-ul REST
NestJS; și (5)~aplicația mobilă React Native.  Integrarea verticală a
acestor componente — firmware \(\leftrightarrow\) broker \(\leftrightarrow\)
\gls{api} \(\leftrightarrow\) aplicație mobilă — asigură un flux de date
continuu, de la citirea brută a senzorilor până la vizualizarea în timp real
și controlul actuatorilor.

Nodul interior, bazat pe modulul ESP32 WROOM-32, integrează doi senzori de
temperatură și umiditate \gls{dht}11, un detector de gaz MQ-2, două
fotorezistențe \gls{ldr} pentru măsurarea luminozității ambientale, trei
\gls{led}-uri, un buzzer, un servo SG90 pentru acționarea draperiilor și un
ventilator de 5\,V.  Camera exterioară ESP32-CAM (AI-Thinker, senzor
OV2640) furnizează flux video \gls{mjpeg} și detectează mișcarea prin
senzorul \gls{pir} conectat la GPIO13.  Toate aceste piese fizice, enumerate
detaliat în \chapref{ch:implementare}, formează stratul primar de percepție
și acțiune al sistemului.

Comunicația intra-rețea se realizează prin protocolul \gls{mqtt}
(\mbox{OASIS}~\cite{oasis2019mqtt}), cu topicuri structurate pe spațiu de
nume \texttt{smarthome/\{nodeId\}/...} și mecanisme de Last~Will pentru
detectarea deconectărilor.  Descoperirea automată a gateway-ului se face prin
\gls{mdns} (\texttt{smarthome.local}), eliminând necesitatea configurării
manuale a adreselor \gls{ip} în aplicația mobilă.

Contribuția originală a lucrării, detaliată în secțiunile dedicate din
\chapref{ch:metodologie} și \chapref{ch:rezultate}, constă în motorul de
detectare a anomaliilor comportamentale, implementat în Python pe
\gls{rpi}.  Motorul extrage 12 trăsături relevante din datele senzorilor,
include variabile temporale și contextuale, normalizează vectorii de
trăsături cu \texttt{StandardScaler} și aplică algoritmul Isolation
Forest~\cite{liu2008isolation}.  Inferența rulează local, în timp real, fără
transmiterea datelor personale în exterior, în acord cu principiile
\emph{local-first}~\cite{kleppmann2019localfirst} și cu cerințele
\gls{gdpr}~\cite{gdpr2016}.

%=============================================================================
\section{Contribuții}
%=============================================================================

Principalele contribuții aduse prin această lucrare sunt:

\begin{enumerate}

\item \textbf{Arhitectură local-first completă pe cinci straturi.}  Sistemul
demonstrează că un ecosistem de casă inteligentă poate fi construit integral
cu componente accesibile (ESP32, Raspberry Pi) și software open-source, fără
a recurge la platforme cloud proprietare.  Această abordare păstrează datele
personale în rețeaua locală și menține funcționalitatea deplină chiar în
absența accesului la internet.

\item \textbf{Motor de detectare a anomaliilor cu Isolation Forest, rulat la
marginea rețelei (\gls{edge}).}  Algoritmul Isolation Forest~\cite{liu2008isolation},
nesupervizat prin natură, se potrivește scenariului \gls{iot} în care
etichetele de anomalie nu sunt disponibile în prealabil.  Implementarea
locală pe \gls{rpi} elimină latența și dependența față de un server remote,
permi\-țând reacția în timp real la evenimente neobișnuite~\cite{shi2016edge,
cook2020anomalyiot}.  Trăsăturile inginerești — incluzând temperatura medie a
celor doi senzori \gls{dht}11, delta de temperatură pe o oră, numărul de
evenimente de mișcare în ultimele 10 minute și indicatori temporali — sporesc
capacitatea de discriminare a modelului față de o abordare bazată exclusiv pe
valorile brute ale senzorilor.

\item \textbf{Pipeline firmware–backend–mobil integrat.}  Firmware-ul
Arduino/C++ al nodurilor publică date la fiecare ciclu de eșantionare;
serviciul \texttt{mqtt.service.ts} din backend consumă aceste date, le
persistă în TimescaleDB și le furnizează aplicației mobile prin Socket.IO,
obținând o latență end-to-end sub o secundă în condiții de rețea locală.

\item \textbf{Interfață mobilă orientată spre siguranță.}  Aplicația React
Native expune un scor de siguranță calculat în timp real, o listă de alerte
cu severitate, istoricul grafic al senzorilor și controlul direct al
actuatorilor, oferind utilizatorului un punct unic de monitorizare și
intervenție.

\end{enumerate}

%=============================================================================
\section{Limitări ale implementării curente}
%=============================================================================

O evaluare onestă a sistemului impune recunoașterea limitărilor prezente:

\paragraph{Model neantrenat pe date reale de durată.}  Modelul Isolation
Forest a fost validat cu date sintetice și cu un set restrâns de înregistrări
reale colectate pe o perioadă scurtă.  Pragul \texttt{ANOMALY\_THRESHOLD =
-0{,}1} și parametrul \texttt{contamination = 0{,}05} au valori
ilustrative, stabilite prin inspecție inițială, nu prin optimizare pe un set
de date reprezentativ acoperind anotimpuri, pattern-uri de ocupare variate și
tipuri diverse de anomalii.  Performanța reală în producție (precizie,
acoperire, scor~F1) rămâne neevaluată riguros și constituie un obiectiv de
cercetare viitoare~\cite{davis2006prcurve,fawcett2006roc}.

\paragraph{Lipsa criptării pe canalul \gls{mqtt}.}  Broker-ul Mosquitto
funcționează pe portul 1883 fără \gls{tls}, ceea ce înseamnă că traficul de
date al senzorilor poate fi interceptat de orice dispozitiv aflat în aceeași
rețea locală.  Deși riscul este limitat la perimetrul \gls{lan}, absența
criptării reprezintă o vulnerabilitate față de atacatorii cu acces fizic sau
wireless la rețea~\cite{alrawi2019sokhome}.

\paragraph{Acoperire redusă a testelor automate.}  Validarea sistemului s-a
bazat preponderent pe teste manuale și pe modulul de simulare a datelor
(\texttt{SIMULATE=true}).  Testele unitare și de integrare acoperă parțial
serviciile backend; nu există teste de regresie automate pentru firmware și
nici teste end-to-end pentru fluxul complet senzor–aplicație mobilă.

\paragraph{Topologie cu un singur nod interior.}  Prototipul curent suportă
un nod interior și o cameră exterioară.  Adăugarea mai multor noduri în
camere diferite necesită extinderea schemei de topicuri \gls{mqtt}, a
modelului de date și a interfeței mobile, fără a fi testată în această
iterație.

%=============================================================================
\section{Direcții de dezvoltare viitoare}
%=============================================================================

Pe baza limitărilor identificate și a potențialului neexplorat al
arhitecturii, sunt propuse următoarele direcții de continuare:

\paragraph{Antrenarea și evaluarea riguroasă a modelului \gls{ml}.}
Colectarea unui set de date reale pe cel puțin 30–90 de zile, acoperind
comportamente normale și evenimente anomale etichetate manual, ar permite
optimizarea hiperparametrilor Isolation Forest și calculul metricilor
cantitative (precizie, recall, AUC-ROC)~\cite{davis2006prcurve,fawcett2006roc,
cook2020anomalyiot}.  Compararea cu alte metode nesupravizate (LOF~\cite{breunig2000lof},
One-Class SVM~\cite{scholkopf2001ocsvm}) ar valida alegerea algoritmică.

\paragraph{Securizarea comunicației \gls{mqtt} cu \gls{tls}.}  Configurarea
Mosquitto cu certificate X.509 și activarea autentificării mutual-TLS pe
portul 8883 ar elimina vulnerabilitatea de interceptare și ar alinia sistemul
la recomandările de securitate pentru dispozitivele \gls{iot}~\cite{alrawi2019sokhome}.

\paragraph{Suport multi-nod.}  Extinderea arhitecturii pentru a suporta mai
multe noduri ESP32 în camere diferite — cu agregarea datelor la nivel de
gateway și vizualizare per-cameră în aplicația mobilă — ar transforma
prototipul actual într-un sistem de monitorizare la scara unei locuințe
complete~\cite{stojkoska2017smarthome}.

\paragraph{Notificări push.}  Integrarea unui serviciu de notificări
(Firebase Cloud Messaging sau un serviciu auto-găzduit) ar permite alertarea
utilizatorului pe dispozitivul mobil chiar și atunci când aplicația este în
fundal sau telefonul nu este în rețeaua locală, crescând semnificativ
caracterul practic al sistemului.

\paragraph{Reguli de automatizare editabile din aplicație.}  Tabelul
\texttt{automation\_rules} există deja în schema bazei de date, iar \gls{api}-ul
expune endpointul \texttt{/api/rules}.  Implementarea unui editor vizual de
reguli în aplicația mobilă — de tipul „dacă temperatura depășește 30°C și
ora este între 10:00 și 22:00, pornește ventilatorul" — ar adăuga un nivel
semnificativ de personalizare, fără modificări de infrastructură.

\paragraph{Optimizări energetice.}  Adoptarea modului deep-sleep pe nodurile
ESP32 și trimiterea datelor la intervale configurabile (în loc de ciclu
continuu) ar reduce consumul energetic al nodurilor alimentate din baterie,
extinzând aplicabilitatea sistemului la scenarii de monitorizare în zone fără
priză permanentă~\cite{shi2016edge,premsankar2018edgeiot}.

\paragraph{Îmbunătățiri ale modelului \gls{ml} prin calcul la margine.}
Pe termen lung, migrarea motorului de inferență pe un accelerator dedicat
(de exemplu, Coral USB, Raspberry Pi AI Kit) sau explorarea unor modele mai
ușoare bazate pe rețele neuronale (autoencoder LSTM pentru serii de timp)
ar putea crește acuratețea detecției și reduce latența de inferență față de
implementarea actuală pe CPU~\cite{chalapathy2019deepanomaly,shi2016edge}.

%=============================================================================
\section{Remarci finale}
%=============================================================================

Sistemul realizat demonstrează viabilitatea unui ecosistem de casă
inteligentă construit integral pe principii \emph{local-first}, cu hardware
accesibil și componente software open-source.  Detecția de anomalii
comportamentale prin Isolation Forest, rulată local pe Raspberry Pi,
reprezintă un pas concret spre sisteme \gls{iot} rezidențiale care protejează
confidențialitatea utilizatorilor fără a sacrifica funcționalitatea.
Limitările identificate — model nevalidat pe date de durată, absența
criptării \gls{mqtt}, acoperire redusă a testelor automate — trasează o
agendă clară de îmbunătățire pentru iterațiile viitoare ale proiectului.


% ── Bibliografie ─────────────────────────────────────────────────────
\printbibliography[heading=bibintoc]

\end{document}
